\documentclass[../DoAn.tex]{subfiles}
\begin{document}

Chương này trình bày quá trình từ thiết kế kiến trúc, xây dựng sản phẩm cho đến triển khai và kiểm thử.
Phần đầu đề cập đến lựa chọn kiến trúc và thiết kế chi tiết các gói thành phần; phần giữa tập trung vào kết quả xây dựng với các thành phần kỹ thuật trọng yếu; phần cuối trình bày kết quả kiểm thử hộp đen và hướng dẫn triển khai hệ thống.

\section{Thiết kế kiến trúc}
\label{section:4.1}

\subsection{Lựa chọn kiến trúc phần mềm}
\label{subsection:4.1.1}

Hệ thống áp dụng \textbf{kiến trúc phân tầng (Layered Architecture)} \cite{fowler2002poeaa,bass2021architecture} gồm bốn tầng từ trên xuống: \textbf{tầng trình bày} (Presentation) $\rightarrow$ \textbf{tầng nghiệp vụ} (Business) $\rightarrow$ \textbf{tầng truy cập dữ liệu} (Persistence) $\rightarrow$ \textbf{tầng dữ liệu} (Database).
Trong cài đặt cụ thể, tầng trình bày được tổ chức nội bộ thành Route, Middleware và Controller; tầng nghiệp vụ là các Service; tầng truy cập dữ liệu gồm Repository và Prisma Client; tầng dữ liệu là PostgreSQL.
Mỗi tầng chỉ tương tác với tầng liền kề, không truy cập vượt tầng (Hình~\ref{fig:kientruc-phantang}).
Đây \textbf{không phải mô hình MVC thuần}.
Backend là API thuần, không có tầng View; phần giao diện do frontend Next.js đảm nhiệm như một tầng client tách biệt.
Hệ thống bổ sung tầng Service cho logic nghiệp vụ và tầng Repository cho truy cập dữ liệu, giúp hai phần này tách bạch rõ ràng.

Lớp Repository được tách khỏi Service để phân định rạch ròi logic nghiệp vụ với truy cập dữ liệu.
Ban đầu, tầng Service gọi thẳng Prisma Client để truy vấn, khiến logic nghiệp vụ phụ thuộc trực tiếp vào ORM.
Sau khi tách, Service chỉ làm việc qua các phương thức của Repository, nhờ đó có thể thay đổi hoặc thay thế ORM mà không ảnh hưởng tới logic nghiệp vụ.

Các kiến trúc khác đã được cân nhắc nhưng không được chọn:

\begin{itemize}
    \item \textbf{MVC thuần} không quy định rõ tầng Service và Repository. Với một dự án có 23 model và nhiều quan hệ nghiệp vụ phức tạp như PM QLKT, nếu chỉ tổ chức theo MVC đơn giản thì logic nghiệp vụ và truy cập dữ liệu dễ bị trộn lẫn trong Controller hoặc Service, gây khó bảo trì và kiểm thử.
    \item \textbf{Clean Architecture / Hexagonal} đầy đủ với Use Case Layer, Entity Layer, Adapter Layer quá nặng so với quy mô dự án, vì lớp Use Case sẽ gần như trùng với tầng Service sẵn có.
    \item \textbf{Microservices} không phù hợp với môi trường mạng nội bộ một máy chủ; chia tách dịch vụ chỉ làm tăng độ phức tạp triển khai mà không mang lại lợi ích về khả năng chịu tải.
\end{itemize}

\begin{figure}[H]
    \centering
    \includegraphics[width=\textwidth]{KienTrucPhanTang}
    \caption{Sơ đồ kiến trúc phân tầng (Layered Architecture) bốn tầng phía backend}
    \label{fig:kientruc-phantang}
\end{figure}

\subsection{Tổng quan hệ thống}
\label{subsection:4.1.2}

Sản phẩm được tổ chức thành ba thành phần độc lập triển khai:

\begin{itemize}
    \item \textbf{Frontend (FE-QLKT)}: ứng dụng Next.js 14 chạy bằng Node.js. Cung cấp giao diện cho cả bốn vai trò người dùng. Kết nối tới backend qua REST API (HTTP) và Socket.IO (WebSocket).
    \item \textbf{Backend (BE-QLKT)}: API server Express + TypeScript chạy bằng Node.js. Cung cấp REST endpoint, Socket.IO server, thực hiện toàn bộ logic nghiệp vụ và quản lý phiên đăng nhập.
    \item \textbf{Cơ sở dữ liệu}: PostgreSQL chạy trong tiến trình riêng. Backend kết nối qua Prisma Client, không cho phép FE truy cập trực tiếp.
\end{itemize}

Cả ba thành phần này được triển khai trên cùng một máy chủ Windows Server của Học viện và giao tiếp qua localhost.
Bên ngoài, người dùng truy cập qua trình duyệt thông qua mạng LAN; FE phục vụ giao diện, còn trình duyệt gọi trực tiếp tới BE qua REST API và Socket.IO.

\subsection{Thiết kế tổng quan các gói}
\label{subsection:4.1.3}

Cấu trúc thư mục mã nguồn được tổ chức theo nguyên tắc đơn trách nhiệm và đặt tên theo tiếng Anh chuẩn để dễ tham chiếu trong các sơ đồ.

\subsubsection{Gói phía frontend}

Mã nguồn frontend được tổ chức trong thư mục nguồn của module \texttt{FE-QLKT}.
Gói \texttt{app/} chứa cây định tuyến của Next.js App Router với bốn nhánh tương ứng bốn vai trò người dùng.
Gói \texttt{components/} chứa các thành phần React dùng chung, tách thành nhiều thư mục con theo miền nghiệp vụ như xác thực, đề xuất, quân nhân và nhật ký hệ thống.
Gói \texttt{contexts/} giữ trạng thái đăng nhập dùng chung cho toàn ứng dụng, còn gói \texttt{hooks/} tập hợp các React hook tự viết phục vụ gọi API, kiểm tra phân quyền và kết nối thời gian thực.
Gói \texttt{lib/} gom các module phụ trợ: lớp gọi REST tới backend (chia theo miền nghiệp vụ), các hàm hỗ trợ hiển thị chuỗi danh hiệu và các schema Zod để xác thực biểu mẫu.
Cuối cùng, gói \texttt{constants/} chứa các hằng số dùng chung như danh sách vai trò, trạng thái và danh hiệu, còn gói \texttt{configs/} giữ cấu hình môi trường của ứng dụng.
Cấu trúc gói phía frontend được thể hiện ở Hình~\ref{fig:goi-frontend}.

\begin{figure}[H]
    \centering
    \includegraphics[width=\textwidth,height=0.9\textheight,keepaspectratio]{goi-fe}
    \caption{Biểu đồ gói phía frontend}
    \label{fig:goi-frontend}
\end{figure}

\subsubsection{Gói phía backend}

Mã nguồn backend được tổ chức trong thư mục nguồn của module \texttt{BE-QLKT} theo các tầng kiến trúc đã trình bày ở mục~\ref{subsection:4.1.1}.
Gói \texttt{routes/} định nghĩa các điểm cuối REST API, mỗi miền nghiệp vụ một tệp.
Gói \texttt{middlewares/} tập hợp các middleware dùng chung cho xác thực và phân quyền, ghi nhật ký và kiểm tra dữ liệu bằng Zod.
Gói \texttt{controllers/} điều hướng yêu cầu HTTP từ tầng định tuyến xuống tầng nghiệp vụ và được giữ mỏng.
Gói \texttt{services/} chứa logic nghiệp vụ chính, tách tiếp thành các thư mục con theo từng mảng: registry các chiến lược của bảy loại đề xuất, module xét điều kiện chuỗi dùng chung cho cá nhân và đơn vị, cùng các module xử lý chuỗi danh hiệu của cá nhân và của đơn vị.
Các service nghiệp vụ lớn được tách theo mẫu facade thành thư mục con, mỗi mối quan tâm một tệp như phần nhập Excel hay phần khai báo kiểu. Bốn service \texttt{account}, \texttt{commemorativeMedal}, \texttt{militaryFlag} và \texttt{scientificAchievement} được tổ chức theo mẫu này.
Gói \texttt{repositories/} đóng gói toàn bộ truy cập Prisma, gồm 21 repository, mỗi model nghiệp vụ một tệp.
Gói \texttt{helpers/} gom các tiện ích phụ trợ theo nhóm như nhật ký, thông báo, Excel và kiểm tra điều kiện khen thưởng; việc lọc dữ liệu theo phạm vi đơn vị là một tiện ích trong \texttt{controllerHelper}. Gói \texttt{validations/} chứa các schema Zod cho từng route.
Riêng tệp khai báo schema ở cấp gốc của module backend định nghĩa toàn bộ 23 model.
Cấu trúc gói phía backend được thể hiện ở Hình~\ref{fig:goi-backend}.

\begin{figure}[H]
    \centering
    \includegraphics[width=\textwidth,height=0.9\textheight,keepaspectratio]{goi-be}
    \caption{Biểu đồ gói phía backend}
    \label{fig:goi-backend}
\end{figure}

\subsection{Thiết kế chi tiết gói nghiệp vụ khen thưởng cá nhân hằng năm}
\label{subsection:4.1.4}

Để minh hoạ cách các tầng phối hợp trong một module nghiệp vụ cụ thể, phần này phân tích module chuỗi danh hiệu hằng năm cá nhân, module phức tạp nhất hệ thống.
Biểu đồ chi tiết các gói thành phần và quan hệ phụ thuộc của module được trình bày ở Hình~\ref{fig:goi-chi-tiet-hangnam}.

\begin{figure}[H]
    \centering
    \includegraphics[width=\textwidth,height=0.9\textheight,keepaspectratio]{so-do-goi-cnhn}
    \caption{Biểu đồ chi tiết gói nghiệp vụ khen thưởng cá nhân hằng năm}
    \label{fig:goi-chi-tiet-hangnam}
\end{figure}

Khi người dùng truy cập trang hồ sơ một quân nhân, dòng dữ liệu đi qua các tầng theo thứ tự:

\begin{enumerate}
    \item \textbf{Route}: tiếp nhận yêu cầu lấy hồ sơ theo mã quân nhân và cho đi qua bước xác thực token; quyền xem hồ sơ của chính mình được kiểm ở tầng nghiệp vụ nên route này không cần kiểm tra vai trò riêng.
    \item \textbf{Controller}: đọc mã quân nhân từ yêu cầu, gọi service tương ứng rồi đóng gói kết quả trả về theo định dạng phản hồi chuẩn.
    \item \textbf{Service}: nạp danh sách danh hiệu hằng năm của quân nhân qua repository, tính ngữ cảnh chuỗi, lần lượt xét từng cấp danh hiệu rồi ghi kết quả vào bảng hồ sơ điều kiện \texttt{HoSoHangNam}.
    \item \textbf{Module xét điều kiện}: chứa hàm thuần thực thi quy tắc chuỗi; vì không truy vấn cơ sở dữ liệu nên hàm này dễ kiểm chứng với mọi bộ đầu vào.
    \item \textbf{Repository}: cung cấp các phương thức đọc, ghi dữ liệu danh hiệu và che giấu chi tiết truy cập Prisma khỏi tầng nghiệp vụ.
    \item \textbf{Prisma Client}: thực thi truy vấn SQL tới bảng \texttt{danh\_hieu\_hang\_nam}.
\end{enumerate}

Việc tách tầng Service (điều phối nghiệp vụ) khỏi module xét điều kiện (quy tắc thuần) giải quyết tốt mối quan tâm về tính kiểm chứng và bảo trì.
Phần quy tắc hoàn toàn thuần nên kiểm chứng được với mọi kịch bản đầu vào mà không cần cơ sở dữ liệu, trong khi tầng Service chỉ lo nạp dữ liệu, gọi quy tắc và ghi kết quả; ranh giới giữa hai phần vì thế rõ ràng và dễ bảo trì.

\subsection{Thiết kế chi tiết gói nghiệp vụ khen thưởng đơn vị hằng năm}
\label{subsection:goi-chi-tiet-donvi}

Gói khen thưởng đơn vị hằng năm có cấu trúc tầng tương tự gói cá nhân, được trình bày ở Hình~\ref{fig:goi-chi-tiet-donvi}.
Điểm khác biệt nằm ở phân quyền và dữ liệu nghiệp vụ, không nằm ở kiến trúc.

\begin{figure}[H]
    \centering
    \includegraphics[width=\textwidth,height=0.9\textheight,keepaspectratio]{so-do-goi-dvhn}
    \caption{Biểu đồ chi tiết gói nghiệp vụ khen thưởng đơn vị hằng năm}
    \label{fig:goi-chi-tiet-donvi}
\end{figure}

Dòng dữ liệu khi tính lại điều kiện khen thưởng cho một đơn vị đi qua các tầng:

\begin{enumerate}
    \item \textbf{Route}: vì chỉ Cán bộ Phòng Chính trị và Chỉ huy đơn vị được thao tác trên khen thưởng đơn vị, khác với gói cá nhân, route này chặn vai trò ngay từ đầu (xác thực token rồi kiểm tra quyền).
    \item \textbf{Controller}: chuyển yêu cầu xuống tầng nghiệp vụ.
    \item \textbf{Service}: đóng vai trò lớp điều phối, uỷ thác phần xét điều kiện cho một module riêng; nó nạp danh sách danh hiệu đơn vị qua repository, đếm số năm Đơn vị Quyết thắng liên tục, rồi gọi quy tắc chuỗi và ghi kết quả vào bảng hồ sơ điều kiện \texttt{HoSoDonViHangNam}.
    \item \textbf{Module xét điều kiện}: dùng lại đúng hàm xét chuỗi chung, chỉ thay bằng bộ cấu hình dành cho đơn vị.
    \item \textbf{Repository}: đóng gói truy cập Prisma tới hai bảng danh hiệu đơn vị và hồ sơ đơn vị.
\end{enumerate}

Điểm thiết kế đáng chú ý là cả hai gói (cá nhân và đơn vị) dùng chung một hàm xét chuỗi thuần, chỉ khác nhau ở bộ cấu hình.
Nhờ vậy, quy tắc xét chuỗi danh hiệu chỉ được mô tả ở một nơi duy nhất, hai luồng không thể lệch nhau, và việc bổ sung danh hiệu mới chỉ cần thêm cấu hình.
Khác biệt nghiệp vụ thể hiện rõ trên biểu đồ: gói cá nhân có thêm dữ liệu thành tích khoa học, yêu cầu nghiên cứu khoa học mỗi năm và xét theo từng quân nhân, trong khi gói đơn vị không xét nghiên cứu khoa học.

\section{Thiết kế chi tiết}
\label{section:4.2}

\subsection{Thiết kế giao diện}
\label{subsection:4.2.1}

Giao diện hệ thống được phác thảo bằng wireframe trước khi triển khai bằng React + Ant Design.
Bảy wireframe dưới đây thể hiện bố cục các trang trọng tâm, từ màn hình đăng nhập đến các trang nghiệp vụ chính.

\textbf{Màn hình đăng nhập} (Hình~\ref{fig:wf-login}) bố trí một thẻ đăng nhập căn giữa trên nền ảnh, gồm logo, tiêu đề, hai trường tên đăng nhập và mật khẩu cùng nút đăng nhập.
Đây là điểm vào duy nhất của hệ thống.
Mọi vai trò đều xác thực qua màn hình này trước khi được điều hướng tới trang chủ tương ứng.

\begin{figure}[H]
    \centering
    \includegraphics[width=0.7\textwidth]{wireframe/wireframe-login}
    \caption{Wireframe màn hình đăng nhập}
    \label{fig:wf-login}
\end{figure}

\textbf{Trang chủ Cán bộ Phòng Chính trị} (Hình~\ref{fig:wf-dashboard}) có cấu trúc ba khu vực.
Thanh điều hướng bên trái là cây cơ quan đơn vị và đơn vị trực thuộc có thể thu gọn.
Khu vực dashboard ở giữa gồm bốn thẻ thống kê là tổng quân nhân, tổng đơn vị, đề xuất chờ duyệt và khen thưởng tháng này, kèm các biểu đồ phân bố.
Khu vực phụ bên phải hiển thị danh sách thông báo thời gian thực.
Bố cục ba khu vực này được giữ nhất quán xuyên suốt các trang nghiệp vụ chính như Quản lý quân nhân, Quản lý đề xuất và Quản lý quyết định để giảm thời gian làm quen cho người dùng mới.

\begin{figure}[H]
    \centering
    \includegraphics[width=0.85\textwidth]{wireframe/wireframe-dashboard}
    \caption{Wireframe trang chủ Cán bộ Phòng Chính trị với các thẻ thống kê và biểu đồ}
    \label{fig:wf-dashboard}
\end{figure}

\textbf{Trang danh sách quân nhân} (Hình~\ref{fig:wf-personnel-list}) gồm thanh điều hướng bên trái, các thẻ thống kê nhanh, thanh tìm kiếm và bộ lọc nhiều tiêu chí ở đầu, cùng bảng dữ liệu phân trang có cột thao tác cuối dòng.
Mẫu bố cục ``bộ lọc trên, bảng dưới, thao tác cuối dòng'' này được tái sử dụng cho hầu hết các trang danh sách (quân nhân, đề xuất, quyết định, khen thưởng).

\begin{figure}[H]
    \centering
    \includegraphics[width=0.85\textwidth]{wireframe/wireframe-personnel-list}
    \caption{Wireframe trang danh sách quân nhân}
    \label{fig:wf-personnel-list}
\end{figure}

\textbf{Trang chi tiết quân nhân} (Hình~\ref{fig:wf-personnel-detail}) hiển thị thông tin định danh ở đầu (ảnh đại diện, họ tên, mã và vai trò) và chia phần thân thành các thẻ tương ứng các nhóm khen thưởng: Hằng năm, Niên hạn, Cống hiến, Quân kỳ, Kỷ niệm chương, NCKH, Đột xuất.
Mỗi thẻ có bảng dữ liệu kèm thanh tiến trình hiển thị các chu kỳ đang đến mốc đề nghị (ví dụ ``BKBQP: 1/2 năm CSTĐCS đã tích luỹ'').
Cách tiếp cận trực quan này giúp Chỉ huy đơn vị nhanh chóng nhận biết quân nhân nào sắp đủ điều kiện.

\begin{figure}[H]
    \centering
    \includegraphics[width=0.85\textwidth]{wireframe/wireframe-personnel-detail}
    \caption{Wireframe trang chi tiết quân nhân với các thẻ khen thưởng}
    \label{fig:wf-personnel-detail}
\end{figure}

\textbf{Trang danh sách đề xuất} (Hình~\ref{fig:wf-proposal-list}) bổ sung hàng thẻ trạng thái (Tất cả, Chờ duyệt, Đã duyệt, Từ chối) phía trên bảng, mỗi dòng kèm nhãn trạng thái màu và nút xem chi tiết.
Nút ``Tạo mới'' ở góc phải dẫn tới biểu mẫu tạo đề xuất.

\begin{figure}[H]
    \centering
    \includegraphics[width=0.85\textwidth]{wireframe/wireframe-proposal-list}
    \caption{Wireframe trang danh sách đề xuất khen thưởng}
    \label{fig:wf-proposal-list}
\end{figure}

\textbf{Biểu mẫu tạo đề xuất} (Hình~\ref{fig:wf-proposal-create}) áp dụng biểu mẫu năm bước (multi-step wizard) với thanh stepper ở đầu.
Bước 1: chọn loại đề xuất (radio button có icon đại diện cho từng nhóm) và năm/tháng.
Bước 2: chọn quân nhân hoặc đơn vị từ bảng có lọc theo nhiều tiêu chí.
Bước 3: thiết lập danh hiệu cụ thể cho từng đối tượng đã chọn, kèm gợi ý từ hệ thống.
Bước 4: tải lên tệp đính kèm minh chứng.
Bước 5: xem lại toàn bộ thông tin và gửi đề xuất.
Người dùng có thể quay lại các bước trước mà không mất dữ liệu nhờ trạng thái được lưu vào \texttt{sessionStorage} của trình duyệt.

\begin{figure}[H]
    \centering
    \includegraphics[width=0.85\textwidth]{wireframe/wireframe-proposal-create}
    \caption{Wireframe biểu mẫu tạo đề xuất năm bước}
    \label{fig:wf-proposal-create}
\end{figure}

\textbf{Trang danh sách khen thưởng} (Hình~\ref{fig:wf-awards}) dùng hàng thẻ chọn loại danh hiệu phía trên, kèm thanh lọc và bảng dữ liệu bên dưới.
Cùng một bố cục phục vụ cả bảy loại khen thưởng, chỉ thay đổi các cột hiển thị theo loại đang chọn.

\begin{figure}[H]
    \centering
    \includegraphics[width=0.85\textwidth]{wireframe/wireframe-awards}
    \caption{Wireframe trang danh sách khen thưởng theo loại}
    \label{fig:wf-awards}
\end{figure}

\textbf{Trang nhật ký hệ thống} (Hình~\ref{fig:wf-system-logs}) đặt bảng lọc ở đầu, gồm các thẻ lọc nhanh, ô tìm kiếm, khoảng thời gian và bộ lọc theo nhóm chức năng và hành động; tiếp đến là hàng thẻ thống kê và bảng nhật ký có tiêu đề cột sắp xếp được.
Mỗi dòng có thể mở rộng để xem chi tiết dữ liệu trước và sau khi thay đổi.
Trang này áp dụng đúng mẫu bố cục ``bộ lọc trên, bảng dưới'' đã dùng cho các trang danh sách khác, giúp người dùng truy vết thao tác mà không phải làm quen giao diện mới.

\begin{figure}[H]
    \centering
    \includegraphics[width=0.85\textwidth]{wireframe/wireframe-system-logs}
    \caption{Wireframe trang nhật ký hệ thống}
    \label{fig:wf-system-logs}
\end{figure}

Ngoài tám trang trên, trang nhập Excel (vùng kéo thả tệp kèm bảng xem trước hai thẻ ``Hợp lệ'' / ``Lỗi'', thẻ Lỗi tô đỏ và liệt kê chỉ số dòng) cũng kế thừa cùng khung bố cục ba khu vực nói trên.

\subsection{Hệ thống thiết kế (Design System)}
\label{subsection:4.2.2}

Bảng màu của hệ thống kế thừa mặc định của Ant Design 5, với gam xanh đại dương cho hành động chính, đỏ cho hành động xoá, vàng cho cảnh báo và xanh lá cho thành công.
Một số tinh chỉnh nhỏ qua Tailwind CSS giúp giao diện phù hợp đặc thù môi trường quân sự, trong đó các màu phụ được giảm độ bão hoà để trông trang trọng hơn.

Hệ thống phông chữ chủ đạo dùng phông sans-serif \texttt{Inter} cho phần thân nội dung do khả năng hiển thị tốt với tiếng Việt có dấu.
Khu vực hiển thị số quyết định và mã định danh kỹ thuật dùng phông \texttt{Source Code Pro} để dễ đọc các chuỗi số dài.

Khoảng cách (spacing) được chuẩn hoá theo thang 4 px của Tailwind.
Tất cả các thẻ (Card) có lề trong tối thiểu 16 px để tránh nội dung sát mép.
Bán kính bo góc được chọn 6 px cho mọi phần tử để tạo cảm giác hiện đại nhưng vẫn nghiêm túc.

\subsection{Thiết kế lớp}
\label{subsection:4.2.3}

Phần này trình bày biểu đồ lớp cho module đề xuất khen thưởng (Hình~\ref{fig:lop-dexuat}).
Đây là module nghiệp vụ phức tạp và tiêu biểu nhất của hệ thống, đồng thời minh hoạ rõ một quyết định thiết kế đáng chú ý là việc áp dụng mẫu thiết kế Strategy.
Biểu đồ chỉ ra các lớp và interface cốt lõi cùng các phương thức quan trọng và mối quan hệ giữa chúng.

Interface \texttt{ProposalStrategy} áp dụng mẫu thiết kế Strategy \cite{gof1994patterns}, quy định bốn phương thức bắt buộc cho mỗi loại đề xuất: dựng dữ liệu khi gửi, kiểm tra khi phê duyệt, ghi dữ liệu trong giao dịch và sinh thông điệp kết quả.
Registry trung tâm ánh xạ tám loại khen thưởng tới chiến lược tương ứng.
Bảy loại đi qua luồng đề xuất và phê duyệt chuẩn có bảy lớp triển khai là \texttt{caNhanHangNamStrategy}, \texttt{donViHangNamStrategy}, \texttt{hccsvvStrategy}, \texttt{hcqkqtStrategy}, \texttt{kncStrategy}, \texttt{hcbvtqStrategy} và \texttt{nckhStrategy}.
Riêng khen thưởng đột xuất (\texttt{DOT\_XUAT}) cố ý được ánh xạ tới \texttt{null} vì không đi qua luồng đề xuất và phê duyệt chuẩn.
Hai loại chỉ trao một huân chương duy nhất là \texttt{hcqkqtStrategy} (HCQKQT) và \texttt{kncStrategy} (KNC) dùng chung hàm \texttt{importSingleMedal} để tránh trùng lặp.
Lớp \texttt{ProposalService} điều phối luồng tạo và phê duyệt đề xuất: nó chọn chiến lược phù hợp qua registry, mỗi chiến lược ghi dữ liệu qua \texttt{proposalRepository} trong một giao dịch và uỷ thác phần kiểm tra điều kiện cho các hàm thuần trong gói \texttt{eligibility/}.

\begin{figure}[H]
    \centering
    \includegraphics[width=\textwidth,height=0.9\textheight,keepaspectratio]{lop-de-xuat}
    \caption{Biểu đồ lớp module đề xuất khen thưởng (mẫu Strategy)}
    \label{fig:lop-dexuat}
\end{figure}

\subsection{Biểu đồ tuần tự}
\label{subsection:4.2.4}

Hai luồng nghiệp vụ phức tạp và tiêu biểu nhất là tạo đề xuất và phê duyệt khen thưởng.
Cả hai được mô hình hoá bằng biểu đồ tuần tự (sequence diagram), thể hiện tương tác giữa các đối tượng theo thời gian.

\textbf{Luồng tạo đề xuất} (Hình~\ref{fig:tuantu-dexuat}): Chỉ huy đơn vị hoàn thành biểu mẫu năm bước; frontend gửi yêu cầu tạo đề xuất tới backend; yêu cầu lần lượt đi qua chuỗi middleware rồi tới tầng nghiệp vụ.
Tại đây, hệ thống chọn chiến lược ứng với loại đề xuất để dựng dữ liệu, ghi bản ghi đề xuất xuống cơ sở dữ liệu, phát một sự kiện Socket.IO báo cho Cán bộ Phòng Chính trị, rồi trả về đề xuất vừa tạo.

\begin{figure}[H]
    \centering
    \includegraphics[width=\textwidth,height=0.9\textheight,keepaspectratio]{seq-de-xuat}
    \caption{Biểu đồ tuần tự luồng đề xuất khen thưởng}
    \label{fig:tuantu-dexuat}
\end{figure}

\textbf{Luồng phê duyệt đề xuất} (Hình~\ref{fig:tuantu-pheduyet}): Cán bộ Phòng Chính trị mở chi tiết đề xuất và bấm ``Phê duyệt''; frontend gửi yêu cầu phê duyệt tới backend.
Sau chuỗi middleware, tầng nghiệp vụ chạy các bước kiểm tra trước (điều kiện, trùng lặp và số quyết định), gắn quyết định cho từng đối tượng, rồi mở một giao dịch Prisma vừa ghi danh hiệu vừa chuyển trạng thái đề xuất sang đã duyệt.
Giao dịch hoàn tất xong, hệ thống tính lại hồ sơ điều kiện, ghi nhật ký và phát sự kiện Socket.IO thông báo kết quả.

\begin{figure}[H]
    \centering
    \includegraphics[width=\textwidth,height=0.9\textheight,keepaspectratio]{seq-phe-duyet}
    \caption{Biểu đồ tuần tự luồng phê duyệt khen thưởng}
    \label{fig:tuantu-pheduyet}
\end{figure}

\subsection{Thiết kế cơ sở dữ liệu}
\label{subsection:4.2.5}

Cơ sở dữ liệu được mô hình hoá bằng Prisma thành \textbf{23 bảng}, liên kết với nhau qua \textbf{47 khoá ngoại}, kèm 13 ràng buộc duy nhất đơn cột và 8 ràng buộc duy nhất tổ hợp.
Lược đồ được trình bày theo hai mức.
Mức thứ nhất là sơ đồ danh mục liệt kê đầy đủ cấu trúc cột của cả 23 bảng, lược bỏ các đường quan hệ để cho cái nhìn tổng thể về thực thể (Hình~\ref{fig:erd-danhmuc}); các quan hệ khoá ngoại được trình bày chi tiết ở các nhóm dữ liệu tiếp theo.
Mức thứ hai gồm sáu sơ đồ nhóm dữ liệu tách theo nghiệp vụ, mỗi sơ đồ chỉ giữ các bảng và quan hệ thuộc nhóm đó để đường khoá ngoại hiện rõ.
Cách phân rã này theo nguyên tắc mỗi sơ đồ con chỉ chứa khoảng năm đến mười bảng nhằm bảo đảm tính đọc được.

\begin{figure}[p]
    \centering
    \includegraphics[width=\textwidth,height=0.92\textheight,keepaspectratio]{Database/erd-0-tat-ca-bang.png}
    \caption{Tổng quan các bảng trong cơ sở dữ liệu}
    \label{fig:erd-danhmuc}
\end{figure}

Do cơ sở dữ liệu của hệ thống gồm nhiều bảng và nhiều quan hệ khoá ngoại, báo cáo không trình bày toàn bộ quan hệ trong một ERD duy nhất.
Thay vào đó, lược đồ được chia thành các nhóm dữ liệu theo nghiệp vụ để mỗi sơ đồ chỉ thể hiện một phần quan hệ có liên quan chặt chẽ với nhau.

\textbf{Nhóm dữ liệu Tổ chức, Quân nhân và Tài khoản} (Hình~\ref{fig:erd-tochuc}) gồm sáu bảng \texttt{CoQuanDonVi}, \texttt{DonViTrucThuoc}, \texttt{ChucVu}, \texttt{QuanNhan}, \texttt{LichSuChucVu}, \texttt{TaiKhoan} với chín khoá ngoại.
Cơ quan đơn vị là đơn vị cha của đơn vị trực thuộc.
Mỗi quân nhân gắn với một chức vụ hiện tại và tối đa một tài khoản đăng nhập.
Bảng \texttt{LichSuChucVu} ghi lại quá trình giữ chức vụ, phục vụ tính thời gian phục vụ xét HCBVTQ.

\begin{figure}[H]
    \centering
    \includegraphics[width=\textwidth,height=0.9\textheight,keepaspectratio]{Database/erd-1-to-chuc.png}
    \caption{ERD nhóm dữ liệu tổ chức, quân nhân và tài khoản}
    \label{fig:erd-tochuc}
\end{figure}

\textbf{Nhóm dữ liệu đề xuất khen thưởng} (Hình~\ref{fig:erd-dexuat}) lấy bảng \texttt{BangDeXuat} làm trung tâm với bốn khoá ngoại trỏ tới đơn vị nộp đề xuất, tài khoản người đề xuất và tài khoản người duyệt.
Dữ liệu chi tiết của từng loại đề xuất được lưu trong các cột JSON, giúp bảng giữ tính đồng nhất mà không phát sinh nhiều bảng trung gian.

\begin{figure}[H]
    \centering
    \includegraphics[width=\textwidth,height=0.9\textheight,keepaspectratio]{Database/erd-2-de-xuat.png}
    \caption{ERD nhóm dữ liệu đề xuất khen thưởng}
    \label{fig:erd-dexuat}
\end{figure}

\textbf{Nhóm dữ liệu khen thưởng và hồ sơ cá nhân} (Hình~\ref{fig:erd-ktcn}) gồm mười bảng khen thưởng và hồ sơ điều kiện có quan hệ trực tiếp với \texttt{QuanNhan} qua mười hai khoá ngoại.
Ba bảng \texttt{HoSoNienHan}, \texttt{HoSoCongHien}, \texttt{HoSoHangNam} là bảng suy diễn, mỗi quân nhân tương ứng với một bản ghi, lưu kết quả tính điều kiện được hệ thống cập nhật lại sau mỗi thay đổi nguồn.

\begin{figure}[H]
    \centering
    \includegraphics[width=\textwidth,height=0.9\textheight,keepaspectratio]{Database/erd-3-khen-thuong-ca-nhan.png}
    \caption{ERD nhóm dữ liệu khen thưởng và hồ sơ cá nhân}
    \label{fig:erd-ktcn}
\end{figure}

\textbf{Nhóm dữ liệu khen thưởng và hồ sơ đơn vị} (Hình~\ref{fig:erd-ktdv}) gồm \texttt{DanhHieuDonViHangNam} và \texttt{HoSoDonViHangNam} với sáu khoá ngoại.
Mỗi bảng liên kết tới cả cơ quan đơn vị lẫn đơn vị trực thuộc tuỳ cấp của đối tượng được khen; riêng \texttt{DanhHieuDonViHangNam} còn tham chiếu tới tài khoản người tạo và tài khoản người duyệt đề nghị.

\begin{figure}[H]
    \centering
    \includegraphics[width=\textwidth,height=0.9\textheight,keepaspectratio]{Database/erd-4-don-vi.png}
    \caption{ERD nhóm dữ liệu khen thưởng và hồ sơ đơn vị}
    \label{fig:erd-ktdv}
\end{figure}

\textbf{Nhóm dữ liệu Quyết định} (Hình~\ref{fig:erd-qd}) tách riêng bảng \texttt{FileQuyetDinh} làm tâm điểm; tám bảng khen thưởng tham chiếu tới cột \texttt{so\_quyet\_dinh} của bảng này qua mười ba khoá ngoại.
Việc tham chiếu vào số quyết định thay vì vào khoá chính \texttt{id} là chủ ý.
Nhờ ràng buộc cập nhật dây chuyền \texttt{ON UPDATE CASCADE}, khi cán bộ sửa lại số quyết định thì thay đổi tự lan sang mọi bảng khen thưởng liên quan.

\begin{figure}[H]
    \centering
    \includegraphics[width=\textwidth,height=0.9\textheight,keepaspectratio]{Database/erd-5-quyet-dinh.png}
    \caption{ERD nhóm dữ liệu quyết định khen thưởng}
    \label{fig:erd-qd}
\end{figure}

\textbf{Nhóm dữ liệu Hệ thống} (Hình~\ref{fig:erd-hethong}) gồm \texttt{SystemLog}, \texttt{ThongBao} và \texttt{SystemSetting}.
Nhật ký hệ thống ghi lại mọi thao tác; thông báo gửi tới tài khoản người nhận và có thể tham chiếu ngược về bản ghi nhật ký đã sinh ra nó.

\begin{figure}[H]
    \centering
    \includegraphics[width=\textwidth,height=0.9\textheight,keepaspectratio]{Database/erd-6-he-thong.png}
    \caption{ERD nhóm dữ liệu nhật ký, thông báo và cấu hình hệ thống}
    \label{fig:erd-hethong}
\end{figure}

Toàn bộ 23 bảng đều dùng khoá chính \dbid{id} kiểu chuỗi sinh tự động theo chuẩn CUID.
Phần dưới tóm tắt vai trò, các thuộc tính chính và quan hệ khoá ngoại của từng bảng, sắp xếp theo sáu nhóm dữ liệu đã trình bày.

\textbf{Nhóm dữ liệu Tổ chức, Quân nhân và Tài khoản.}
\begin{itemize}
    \item \texttt{CoQuanDonVi}: cơ quan/đơn vị cấp trên. Thuộc tính chính: \dbid{ma_don_vi} (duy nhất), \dbid{ten_don_vi}, \dbid{so_luong}. Quan hệ: được đơn vị trực thuộc, chức vụ, quân nhân và các bảng khen thưởng đơn vị tham chiếu.
    \item \texttt{DonViTrucThuoc}: đơn vị trực thuộc cơ quan. Thuộc tính chính: \dbid{ma_don_vi} (duy nhất), \dbid{ten_don_vi}, \dbid{so_luong}. Quan hệ: thuộc một \texttt{CoQuanDonVi}; được chức vụ, quân nhân và các bảng khen thưởng đơn vị tham chiếu.
    \item \texttt{ChucVu}: chức vụ trong đơn vị. Thuộc tính chính: \dbid{ten_chuc_vu}, \dbid{is_manager}, \dbid{he_so_chuc_vu}. Quan hệ: thuộc cơ quan hoặc đơn vị trực thuộc; được \texttt{QuanNhan} và \texttt{LichSuChucVu} tham chiếu; duy nhất theo cặp đơn vị và tên chức vụ.
    \item \texttt{QuanNhan}: hồ sơ quân nhân. Thuộc tính chính: \dbid{cccd} (duy nhất), \dbid{ho_ten}, \dbid{cap_bac}, \dbid{ngay_nhap_ngu}. Quan hệ: gắn một chức vụ và một đơn vị; là tâm của hầu hết bảng khen thưởng, hồ sơ điều kiện và tài khoản.
    \item \texttt{LichSuChucVu}: lịch sử giữ chức vụ của quân nhân. Thuộc tính chính: \dbid{ngay_bat_dau}, \dbid{ngay_ket_thuc}, \dbid{so_thang}. Quan hệ: tham chiếu \texttt{QuanNhan} và \texttt{ChucVu}; phục vụ tính thời gian công tác.
    \item \texttt{TaiKhoan}: tài khoản đăng nhập. Thuộc tính chính: \dbid{username} (duy nhất), \dbid{password_hash}, \dbid{role}, \dbid{refreshToken}. Quan hệ: gắn tối đa một \texttt{QuanNhan}; được đề xuất (người lập, người duyệt), nhật ký và thông báo tham chiếu.
\end{itemize}

\textbf{Nhóm dữ liệu đề xuất khen thưởng.}
\begin{itemize}
    \item \texttt{BangDeXuat}: đề xuất khen thưởng đang chờ duyệt. Thuộc tính chính: \dbid{loai_de_xuat}, \dbid{nam}, \dbid{status}, cùng bốn cột JSON \dbid{data_danh_hieu}, \dbid{data_thanh_tich}, \dbid{data_nien_han}, \dbid{data_cong_hien}. Quan hệ: tham chiếu đơn vị nộp, tài khoản người đề xuất và tài khoản người duyệt.
\end{itemize}

\textbf{Nhóm dữ liệu khen thưởng và hồ sơ cá nhân.}
\begin{itemize}
    \item \texttt{DanhHieuHangNam}: danh hiệu thi đua hằng năm và các cấp danh hiệu chuỗi. Thuộc tính chính: \dbid{nam}, \dbid{danh_hieu}, các cờ \dbid{nhan_bkbqp}, \dbid{nhan_cstdtq}, \dbid{nhan_bkttcp} kèm số quyết định tương ứng. Quan hệ: tham chiếu \texttt{QuanNhan} và bốn lần tới \texttt{FileQuyetDinh}; duy nhất theo cặp quân nhân và năm.
    \item \texttt{ThanhTichKhoaHoc}: thành tích nghiên cứu khoa học. Thuộc tính chính: \dbid{nam}, \dbid{loai} (đề tài hoặc sáng kiến), \dbid{mo_ta}. Quan hệ: tham chiếu \texttt{QuanNhan} và \texttt{FileQuyetDinh}.
    \item \texttt{KhenThuongHCCSVV}: Huy chương Chiến sĩ vẻ vang theo niên hạn. Thuộc tính chính: \dbid{danh_hieu} (hạng), \dbid{nam}, \dbid{thang}, \dbid{thoi_gian} kiểu JSON. Quan hệ: tham chiếu \texttt{QuanNhan} và \texttt{FileQuyetDinh}; duy nhất theo cặp quân nhân và hạng.
    \item \texttt{KhenThuongHCBVTQ}: Huân chương Bảo vệ Tổ quốc theo cống hiến. Thuộc tính chính: \dbid{danh_hieu}, \dbid{nam}, ba cột thời gian theo nhóm hệ số kiểu JSON. Quan hệ: tham chiếu \texttt{QuanNhan} theo quan hệ một-một và \texttt{FileQuyetDinh}.
    \item \texttt{HuanChuongQuanKyQuyetThang}: Huy chương Quân kỳ quyết thắng. Thuộc tính chính: \dbid{nam}, \dbid{thang}, \dbid{thoi_gian}. Quan hệ: tham chiếu \texttt{QuanNhan} theo quan hệ một-một và \texttt{FileQuyetDinh}.
    \item \texttt{KyNiemChuongVSNXDQDNDVN}: Kỷ niệm chương Vì sự nghiệp xây dựng Quân đội nhân dân Việt Nam. Thuộc tính chính: \dbid{nam}, \dbid{thang}, \dbid{thoi_gian}. Quan hệ: tham chiếu \texttt{QuanNhan} theo quan hệ một-một và \texttt{FileQuyetDinh}.
    \item \texttt{KhenThuongDotXuat}: khen thưởng đột xuất cho cá nhân hoặc tập thể. Thuộc tính chính: \dbid{doi_tuong} (cá nhân hoặc tập thể), \dbid{hinh_thuc_khen_thuong}, \dbid{nam}, \dbid{files_dinh_kem}. Quan hệ: tham chiếu \texttt{QuanNhan} hoặc đơn vị, và \texttt{FileQuyetDinh}.
    \item \texttt{HoSoNienHan}: hồ sơ điều kiện HCCSVV, dạng suy diễn. Thuộc tính chính: ba cặp trạng thái và ngày đạt theo từng hạng, \dbid{goi_y}. Quan hệ: tham chiếu \texttt{QuanNhan} theo quan hệ một-một.
    \item \texttt{HoSoCongHien}: hồ sơ điều kiện HCBVTQ, dạng suy diễn. Thuộc tính chính: \dbid{hcbvtq_total_months}, các cột tháng theo nhóm hệ số, trạng thái theo từng hạng. Quan hệ: tham chiếu \texttt{QuanNhan} theo quan hệ một-một.
    \item \texttt{HoSoHangNam}: hồ sơ điều kiện chuỗi danh hiệu hằng năm, dạng suy diễn. Thuộc tính chính: các bộ đếm liên tục \dbid{cstdcs_lien_tuc}, \dbid{bkbqp_lien_tuc}, \dbid{cstdtq_lien_tuc} và các cờ đủ điều kiện đề nghị, \dbid{goi_y}. Quan hệ: tham chiếu \texttt{QuanNhan} theo quan hệ một-một.
\end{itemize}

\textbf{Nhóm dữ liệu khen thưởng và hồ sơ đơn vị.}
\begin{itemize}
    \item \texttt{DanhHieuDonViHangNam}: danh hiệu thi đua của đơn vị. Thuộc tính chính: \dbid{nam}, \dbid{danh_hieu}, các cờ \dbid{nhan_bkbqp}, \dbid{nhan_bkttcp}, \dbid{status}, \dbid{nguoi_tao_id}. Quan hệ: tham chiếu đơn vị và ba lần tới \texttt{FileQuyetDinh}; duy nhất theo cặp đơn vị và năm.
    \item \texttt{HoSoDonViHangNam}: hồ sơ điều kiện khen thưởng đơn vị, dạng suy diễn. Thuộc tính chính: \dbid{tong_dvqt}, \dbid{dvqt_lien_tuc}, các cờ \dbid{du_dieu_kien_bkbqp}, \dbid{du_dieu_kien_bkttcp}. Quan hệ: tham chiếu đơn vị; duy nhất theo cặp đơn vị và năm.
\end{itemize}

\textbf{Nhóm dữ liệu Quyết định.}
\begin{itemize}
    \item \texttt{FileQuyetDinh}: hồ sơ quyết định khen thưởng. Thuộc tính chính: \dbid{so_quyet_dinh} (duy nhất), \dbid{nam}, \dbid{ngay_ky}, \dbid{nguoi_ky}, \dbid{file_path}. Quan hệ: được tám bảng khen thưởng tham chiếu qua cột \dbid{so_quyet_dinh}, tổng cộng 13 khoá ngoại.
\end{itemize}

\textbf{Nhóm dữ liệu Hệ thống.}
\begin{itemize}
    \item \texttt{SystemLog}: nhật ký kiểm toán mọi thao tác. Thuộc tính chính: \dbid{actor_role}, \dbid{action}, \dbid{resource}, \dbid{payload} kiểu JSON, \dbid{ip_address}. Quan hệ: tham chiếu \texttt{TaiKhoan}; được \texttt{ThongBao} tham chiếu.
    \item \texttt{ThongBao}: thông báo thời gian thực. Thuộc tính chính: \dbid{type}, \dbid{title}, \dbid{message}, \dbid{is_read}, \dbid{recipient_role}. Quan hệ: tham chiếu \texttt{TaiKhoan} (người nhận) và \texttt{SystemLog}.
    \item \texttt{SystemSetting}: cấu hình hệ thống dạng khóa-giá trị. Thuộc tính chính: \dbid{key} (duy nhất), \dbid{value}. Quan hệ: độc lập, không có khoá ngoại.
\end{itemize}

Sáu bảng mô tả chi tiết schema sau đây được chọn vì là các bảng trọng yếu nhất, sử dụng nhiều cấu trúc dữ liệu phức tạp như JSON và kiểu thời gian.

\begin{table}[H]
\centering
\caption{Schema bảng \texttt{QuanNhan} (Personnel)}
\label{tab:schema-quannhan}
\footnotesize
\renewcommand{\arraystretch}{1.3}
\setlength{\tabcolsep}{4pt}\begin{tabular}{|R{4.4cm}|R{3.2cm}|R{5.1cm}|}
\hline
\textbf{Thuộc tính} & \textbf{Kiểu dữ liệu} & \textbf{Ý nghĩa} \\ \hline
\texttt{id} & String (PK, CUID) & Mã định danh quân nhân \\ \hline
\texttt{cccd} & VarChar(20) & Số căn cước công dân \\ \hline
\dbid{ho_ten} & VarChar(255) & Họ và tên đầy đủ \\ \hline
\dbid{ngay_sinh} & Date & Ngày sinh \\ \hline
\dbid{gioi_tinh} & VarChar(10) & Giới tính \\ \hline
\dbid{ngay_nhap_ngu} & Date & Ngày nhập ngũ \\ \hline
\dbid{ngay_xuat_ngu} & Date & Ngày xuất ngũ (nếu đã rời) \\ \hline
\dbid{cap_bac} & VarChar(50) & Cấp bậc hiện tại \\ \hline
\dbid{co_quan_don_vi_id} & VarChar(30) & Cơ quan đơn vị quản lý \\ \hline
\dbid{don_vi_truc_thuoc_id} & VarChar(30) & Đơn vị trực thuộc \\ \hline
\dbid{chuc_vu_id} & VarChar(30) & Chức vụ hiện giữ \\ \hline
\texttt{createdAt} & Timestamp & Thời điểm tạo bản ghi \\ \hline
\texttt{updatedAt} & Timestamp & Thời điểm cập nhật gần nhất \\ \hline
\end{tabular}
\end{table}

Ràng buộc nghiệp vụ: chỉ một trong \texttt{co\_quan\_don\_vi\_id} và \texttt{don\_vi\_truc\_thuoc\_id} được điền.
Khi cập nhật chuyển đơn vị, cần dùng giao dịch để cập nhật cả bộ đếm \texttt{so\_luong} của hai đơn vị liên quan.

\begin{table}[H]
\centering
\caption{Schema bảng \texttt{BangDeXuat} (Proposal)}
\label{tab:schema-dexuat}
\footnotesize
\renewcommand{\arraystretch}{1.3}
\setlength{\tabcolsep}{4pt}\begin{tabular}{|R{4.4cm}|R{3.2cm}|R{5.1cm}|}
\hline
\textbf{Thuộc tính} & \textbf{Kiểu dữ liệu} & \textbf{Ý nghĩa} \\ \hline
\texttt{id} & String (PK, CUID) & Mã định danh đề xuất \\ \hline
\dbid{loai_de_xuat} & VarChar(20) & Một trong bảy loại đề xuất \\ \hline
\dbid{nguoi_de_xuat_id} & VarChar(30) & Tài khoản người tạo đề xuất \\ \hline
\texttt{nam} & Int & Năm đề nghị \\ \hline
\texttt{thang} & Int & Tháng đề nghị (chỉ một số loại) \\ \hline
\texttt{status} & VarChar(20) & Trạng thái: chờ duyệt, đã duyệt, từ chối \\ \hline
\dbid{data_danh_hieu} & JSON & Dữ liệu chi tiết danh hiệu hằng năm \\ \hline
\dbid{data_thanh_tich} & JSON & Dữ liệu thành tích nghiên cứu khoa học \\ \hline
\dbid{data_nien_han} & JSON & Dữ liệu xét Huy chương Chiến sĩ vẻ vang \\ \hline
\dbid{data_cong_hien} & JSON & Dữ liệu xét Huân chương Bảo vệ Tổ quốc \\ \hline
\dbid{files_attached} & JSON & Danh sách tệp đính kèm \\ \hline
\dbid{rejection_reason} & Text & Lý do từ chối (nếu bị bác) \\ \hline
\dbid{nguoi_duyet_id} & VarChar(30) & Tài khoản người phê duyệt \\ \hline
\dbid{ngay_duyet} & Timestamp & Thời điểm phê duyệt \\ \hline
\end{tabular}
\end{table}

Việc dùng JSON cho 4 trường \texttt{data\_*} là chủ ý.
Mỗi loại đề xuất có lược đồ con riêng nên việc tạo nhiều bảng trung gian sẽ phá vỡ tính đồng nhất của bảng \texttt{BangDeXuat}.
JSON đủ để đảm bảo lược đồ linh hoạt; xác thực được đảm nhiệm ở tầng Zod và Strategy tương ứng từng loại.

\begin{table}[H]
\centering
\caption{Schema bảng \texttt{DanhHieuHangNam} (Annual Award)}
\label{tab:schema-dhhn}
\footnotesize
\renewcommand{\arraystretch}{1.3}
\setlength{\tabcolsep}{4pt}\begin{tabular}{|R{4.4cm}|R{3.2cm}|R{5.1cm}|}
\hline
\textbf{Thuộc tính} & \textbf{Kiểu dữ liệu} & \textbf{Ý nghĩa} \\ \hline
\texttt{id} & String (PK, CUID) & Mã định danh bản ghi \\ \hline
\dbid{quan_nhan_id} & VarChar(30) & Quân nhân được xét \\ \hline
\texttt{nam} & Int & Năm xét danh hiệu \\ \hline
\dbid{danh_hieu} & VarChar(20) & Mã danh hiệu cơ bản (CSTĐCS, CSTT) \\ \hline
\dbid{nhan_bkbqp} & Boolean & Đã nhận BKBQP trong năm \\ \hline
\dbid{so_quyet_dinh_bkbqp} & VarChar(100) & Số quyết định BKBQP \\ \hline
\dbid{nhan_cstdtq} & Boolean & Đã nhận CSTĐTQ trong năm \\ \hline
\dbid{so_quyet_dinh_cstdtq} & VarChar(100) & Số quyết định CSTĐTQ \\ \hline
\dbid{nhan_bkttcp} & Boolean & Đã nhận BKTTCP trong năm \\ \hline
\dbid{so_quyet_dinh_bkttcp} & VarChar(100) & Số quyết định BKTTCP \\ \hline
\dbid{ghi_chu} & Text & Ghi chú bổ sung \\ \hline
\end{tabular}
\end{table}

Ràng buộc duy nhất \texttt{(quan\_nhan\_id, nam)} đảm bảo mỗi quân nhân chỉ có một bản ghi cho một năm.
Các trường \texttt{nhan\_*} cho phép cộng dồn các cấp danh hiệu chuỗi vào cùng năm.

\begin{table}[H]
\centering
\caption{Schema bảng \texttt{LichSuChucVu} (Position History)}
\label{tab:schema-lsc}
\footnotesize
\renewcommand{\arraystretch}{1.3}
\setlength{\tabcolsep}{4pt}\begin{tabular}{|R{4.4cm}|R{3.2cm}|R{5.1cm}|}
\hline
\textbf{Thuộc tính} & \textbf{Kiểu dữ liệu} & \textbf{Ý nghĩa} \\ \hline
\texttt{id} & String (PK, CUID) & Mã định danh bản ghi \\ \hline
\dbid{quan_nhan_id} & VarChar(30) & Quân nhân giữ chức vụ \\ \hline
\dbid{chuc_vu_id} & VarChar(30) & Chức vụ được giữ \\ \hline
\dbid{ngay_bat_dau} & Date & Ngày bắt đầu giữ chức vụ \\ \hline
\dbid{ngay_ket_thuc} & Date & Ngày kết thúc (để trống nếu đang giữ) \\ \hline
\dbid{he_so_chuc_vu} & Float & Hệ số chức vụ tại thời điểm đó \\ \hline
\dbid{so_thang} & Int & Số tháng đã giữ \\ \hline
\end{tabular}
\end{table}

Hàm tính lại số tháng tự động cập nhật trường \texttt{so\_thang} mỗi khi có thay đổi mốc thời gian, phục vụ thuật toán xét HCBVTQ.

\begin{table}[H]
\centering
\caption{Schema bảng \texttt{SystemLog} (Audit Log)}
\label{tab:schema-syslog}
\footnotesize
\renewcommand{\arraystretch}{1.3}
\setlength{\tabcolsep}{4pt}\begin{tabular}{|R{4.4cm}|R{3.2cm}|R{5.1cm}|}
\hline
\textbf{Thuộc tính} & \textbf{Kiểu dữ liệu} & \textbf{Ý nghĩa} \\ \hline
\texttt{id} & String (PK, CUID) & Mã định danh bản ghi \\ \hline
\dbid{nguoi_thuc_hien_id} & VarChar(30) & Người thực hiện (trống nếu là hệ thống) \\ \hline
\dbid{actor_role} & VarChar(20) & Vai trò tại thời điểm thao tác \\ \hline
\texttt{action} & VarChar(50) & Loại hành động: tạo, cập nhật, xoá... \\ \hline
\texttt{resource} & VarChar(50) & Loại tài nguyên bị tác động \\ \hline
\dbid{tai_nguyen_id} & VarChar(30) & Mã tài nguyên bị tác động \\ \hline
\texttt{description} & VarChar(500) & Mô tả hành động \\ \hline
\texttt{payload} & JSON & Dữ liệu trước và sau khi thay đổi \\ \hline
\dbid{ip_address} & VarChar(45) & Địa chỉ IP người thực hiện \\ \hline
\dbid{user_agent} & Text & Trình duyệt người thực hiện \\ \hline
\texttt{createdAt} & Timestamp & Thời điểm ghi nhật ký \\ \hline
\end{tabular}
\end{table}

\begin{table}[H]
\centering
\caption{Schema bảng \texttt{HoSoHangNam} (Annual Profile, Derived)}
\label{tab:schema-hsnn}
\footnotesize
\renewcommand{\arraystretch}{1.3}
\setlength{\tabcolsep}{4pt}\begin{tabular}{|R{4.4cm}|R{3.2cm}|R{5.1cm}|}
\hline
\textbf{Thuộc tính} & \textbf{Kiểu dữ liệu} & \textbf{Ý nghĩa} \\ \hline
\texttt{id} & String (PK, CUID) & Mã định danh hồ sơ \\ \hline
\dbid{quan_nhan_id} & VarChar(30) & Quân nhân sở hữu hồ sơ \\ \hline
\dbid{cstdcs_lien_tuc} & Int & Số năm CSTĐCS liên tục hiện tại \\ \hline
\dbid{nckh_lien_tuc} & Int & Số năm NCKH liên tục \\ \hline
\dbid{bkbqp_lien_tuc} & Int & Số năm từ lần nhận BKBQP gần nhất \\ \hline
\dbid{cstdtq_lien_tuc} & Int & Số năm từ lần nhận CSTĐTQ gần nhất \\ \hline
\dbid{du_dieu_kien_bkbqp} & Boolean & Đủ điều kiện đề nghị BKBQP \\ \hline
\dbid{du_dieu_kien_cstdtq} & Boolean & Đủ điều kiện đề nghị CSTĐTQ \\ \hline
\dbid{du_dieu_kien_bkttcp} & Boolean & Đủ điều kiện đề nghị BKTTCP \\ \hline
\dbid{goi_y} & Text & Thông điệp gợi ý \\ \hline
\texttt{updatedAt} & Timestamp & Thời điểm tính lại gần nhất \\ \hline
\end{tabular}
\end{table}

\section{Xây dựng ứng dụng}
\label{section:4.3}

\subsection{Thư viện và công cụ sử dụng}
\label{subsection:4.3.1}

Sản phẩm dùng các thư viện mã nguồn mở liệt kê tại Bảng \ref{tab:thuvien}, phân theo ba nhóm: tầng giao diện, tầng máy chủ và các tiện ích vận hành.

\begin{table}[H]
\centering
\caption{Danh sách thư viện và công cụ sử dụng}
\label{tab:thuvien}
\footnotesize
\begin{tabular}{|p{5.3cm}|p{2.8cm}|p{4.6cm}|}
\hline
\textbf{Mục đích} & \textbf{Công cụ} & \textbf{Địa chỉ URL} \\ \hline
Khung phát triển giao diện người dùng & Next.js & \url{https://nextjs.org} \\ \hline
Thư viện giao diện nền tảng & React & \url{https://react.dev} \\ \hline
Thư viện thành phần UI & Ant Design & \url{https://ant.design} \\ \hline
Khung CSS utility & Tailwind CSS & \url{https://tailwindcss.com} \\ \hline
Môi trường chạy phía máy chủ & Node.js & \url{https://nodejs.org} \\ \hline
Khung phát triển phía máy chủ & Express & \url{https://expressjs.com} \\ \hline
Ngôn ngữ kiểu tĩnh & TypeScript & \url{https://www.typescriptlang.org} \\ \hline
Hệ quản trị cơ sở dữ liệu & PostgreSQL & \url{https://www.postgresql.org} \\ \hline
ORM truy cập dữ liệu & Prisma & \url{https://www.prisma.io} \\ \hline
Xác thực JWT & jsonwebtoken & \url{https://github.com/auth0/node-jsonwebtoken} \\ \hline
Băm mật khẩu & bcrypt & \url{https://www.npmjs.com/package/bcrypt} \\ \hline
Xác thực dữ liệu & Zod & \url{https://zod.dev} \\ \hline
Thông báo thời gian thực & Socket.IO & \url{https://socket.io} \\ \hline
Đọc và ghi Excel & ExcelJS & \url{https://github.com/exceljs/exceljs} \\ \hline
Lập lịch nội bộ & node-cron & \url{https://www.npmjs.com/package/node-cron} \\ \hline
Quản lý tiến trình production & PM2 & \url{https://pm2.keymetrics.io} \\ \hline
Header bảo mật HTTP & Helmet & \url{https://helmetjs.github.io} \\ \hline
Cấu hình CORS & cors & \url{https://www.npmjs.com/package/cors} \\ \hline
Giới hạn tốc độ truy cập & express-rate-limit & \url{https://www.npmjs.com/package/express-rate-limit} \\ \hline
Quản lý tệp tải lên & Multer & \url{https://github.com/expressjs/multer} \\ \hline
\end{tabular}
\end{table}

Môi trường phát triển dùng Visual Studio Code với extension Prisma, ESLint và Prettier; mã nguồn quản lý qua Git và GitHub; cơ sở dữ liệu được tra cứu trực quan qua Prisma Studio (\texttt{npx prisma studio}).

\subsection{Kết quả đạt được}
\label{subsection:4.3.2}

Sau quá trình xây dựng, hệ thống đã hoàn thiện các chức năng chính sau:

\begin{itemize}
    \item Quản lý quân nhân: thêm, sửa, xoá thủ công qua biểu mẫu (không nhập/xuất Excel); chuyển đơn vị và đổi chức vụ chỉ Cán bộ Phòng Chính trị, giữ nguyên lịch sử khen thưởng; phân quyền theo phạm vi đơn vị cho Chỉ huy đơn vị.
    \item Quản lý đơn vị (cơ quan đơn vị và đơn vị trực thuộc) và chức vụ: cây tổ chức hai cấp; tự cập nhật bộ đếm \texttt{so\_luong} khi quân nhân chuyển đơn vị.
    \item Quản lý lịch sử chức vụ: nhập danh sách giai đoạn giữ chức vụ với hệ số; tự tính \texttt{so\_thang} mỗi khi thay đổi mốc thời gian.
    \item Đề xuất khen thưởng cho cả bảy loại với biểu mẫu năm bước; tự động gợi ý quân nhân đủ điều kiện theo từng loại.
    \item Phê duyệt đề xuất với giao dịch, gắn số quyết định, tải lên PDF, ghi nhật ký đầy đủ trạng thái trước và sau khi thay đổi.
    \item Tính lại điều kiện khen thưởng hằng năm cá nhân và đơn vị tự động sau mỗi thay đổi nguồn; gợi ý đề nghị bằng tiếng Việt.
    \item Nhập Excel hai bước (xem trước và xác nhận với giao dịch) cho cả bảy loại nghiệp vụ; xuất danh sách khen thưởng theo nhiều tiêu chí.
    \item Quản trị tài khoản bốn cấp; đặt lại mật khẩu; nhật ký hệ thống; sao lưu định kỳ; khu vực DevZone cho Quản trị viên.
    \item Thông báo thời gian thực qua Socket.IO khi có đề xuất mới, được phê duyệt, bị từ chối, bị xoá hoặc nhập Excel hoàn tất.
\end{itemize}

\subsection{Minh hoạ các chức năng chính}
\label{subsection:4.3.3}

Phần này trình bày các ảnh chụp màn hình tiêu biểu của sản phẩm trong môi trường phát triển.
Toàn bộ ảnh được chụp trên trình duyệt Chromium ở độ phân giải 1920×1080, sau đó cắt khu vực có nội dung và lưu định dạng PNG.
Các hình được nhóm theo bốn vai trò người dùng (Quản trị viên, Cán bộ Phòng Chính trị, Chỉ huy đơn vị, Người dùng) cùng nhóm màn hình chung; mỗi vai trò có trang chủ (dashboard) riêng với bộ thẻ thống kê và biểu đồ phù hợp phạm vi dữ liệu của vai trò đó.

\subsubsection{Màn hình chung và xác thực}

\textbf{Trang chủ} (Hình~\ref{fig:ui-home}) là trang giới thiệu công khai, nêu tên hệ thống cùng vài số liệu tổng quan và hai lối vào ``Đăng nhập'' / ``Đề xuất khen thưởng''.
\textbf{Màn hình đăng nhập} (Hình~\ref{fig:ui-login}) hiện thực bố cục đã phác ở wireframe (Mục~\ref{subsection:4.2.1}), khác biệt ở phần tinh chỉnh giao diện: thẻ đăng nhập dùng hiệu ứng kính mờ trên nền ảnh Học viện và bổ sung liên kết quên mật khẩu.

\begin{figure}[H]
    \centering
    \includegraphics[width=0.85\textwidth]{UI/01-trang-chu}
    \caption{Trang chủ giới thiệu hệ thống}
    \label{fig:ui-home}
\end{figure}

\begin{figure}[H]
    \centering
    \includegraphics[width=0.85\textwidth]{UI/02-dang-nhap}
    \caption{Màn hình đăng nhập}
    \label{fig:ui-login}
\end{figure}

\subsubsection{Vai trò quản trị viên (\texttt{SuperAdmin})}

\textbf{Trang chủ Quản trị viên} (Hình~\ref{fig:ui-sa-dashboard}) tổng hợp toàn hệ thống với bốn thẻ thống kê (tổng tài khoản, tài khoản mới trong 30 ngày, hoạt động trong 7 ngày, tổng nhật ký) cùng các biểu đồ phân bố vai trò, hoạt động hệ thống và top hành động phổ biến.

\begin{figure}[H]
    \centering
    \includegraphics[width=0.92\textwidth]{UI/03-superadmin-dashboard}
    \caption{Trang chủ vai trò quản trị viên}
    \label{fig:ui-sa-dashboard}
\end{figure}

\textbf{Quản lý tài khoản} (Hình~\ref{fig:ui-sa-account}) liệt kê toàn bộ tài khoản kèm vai trò, quân nhân liên kết, đơn vị và thao tác sửa, xoá.
\textbf{Tạo tài khoản} (Hình~\ref{fig:ui-sa-account-create}) là biểu mẫu chọn vai trò, đặt tên đăng nhập (gợi ý dùng CCCD), mật khẩu (để trống sẽ dùng mật khẩu mặc định) và gắn cơ quan, đơn vị, chức vụ.

\begin{figure}[H]
    \centering
    \includegraphics[width=0.92\textwidth]{UI/04-superadmin-quan-ly-tai-khoan}
    \caption{Quản lý danh sách tài khoản}
    \label{fig:ui-sa-account}
\end{figure}

\begin{figure}[H]
    \centering
    \includegraphics[width=0.78\textwidth]{UI/05-superadmin-tao-tai-khoan}
    \caption{Biểu mẫu tạo tài khoản mới}
    \label{fig:ui-sa-account-create}
\end{figure}

\textbf{Quản lý cơ quan đơn vị} (Hình~\ref{fig:ui-sa-unit}) hiển thị cây tổ chức hai cấp (cơ quan đơn vị và các đơn vị trực thuộc); mỗi đơn vị có thể mở rộng để xem danh sách quân nhân trực thuộc kèm cấp bậc và chức vụ.

\begin{figure}[H]
    \centering
    \includegraphics[width=0.92\textwidth]{UI/06-superadmin-quan-ly-co-quan-don-vi}
    \caption{Quản lý cơ quan đơn vị và chức vụ}
    \label{fig:ui-sa-unit}
\end{figure}

\textbf{Nhật ký hệ thống} ở vai trò Quản trị viên bao trùm toàn hệ thống và là nơi duy nhất hiển thị nhóm nhật ký sao lưu dữ liệu (Hình~\ref{fig:ui-sa-logs}); các vai trò thấp hơn không thấy nhóm nhật ký này.

\begin{figure}[H]
    \centering
    \includegraphics[width=0.92\textwidth]{UI/31-superadmin-nhat-ky-he-thong}
    \caption{Nhật ký hệ thống vai trò quản trị viên, gồm cả nhật ký sao lưu dữ liệu}
    \label{fig:ui-sa-logs}
\end{figure}

\subsubsection{Vai trò Cán bộ Phòng Chính trị (\texttt{Admin})}

\textbf{Trang chủ Cán bộ Phòng Chính trị} (Hình~\ref{fig:ui-admin-dashboard}) trình bày năm thẻ thống kê phạm vi toàn quân số quản lý (tổng quân nhân, đơn vị, chức vụ, đề xuất chờ duyệt, đề xuất bị từ chối) cùng các biểu đồ NCKH và đề xuất theo loại và trạng thái.

\begin{figure}[H]
    \centering
    \includegraphics[width=0.92\textwidth]{UI/27-admin-dashboard}
    \caption{Trang chủ vai trò Cán bộ Phòng Chính trị}
    \label{fig:ui-admin-dashboard}
\end{figure}

\textbf{Chi tiết quân nhân} chia phần thân thành nhiều thẻ.
Thẻ \textit{Thông tin cơ bản} (Hình~\ref{fig:ui-admin-personnel-info}) hiển thị thông tin định danh, thông tin cá nhân và thông tin Đảng.
Thẻ \textit{Hồ sơ khen thưởng} (Hình~\ref{fig:ui-admin-personnel-award}) trình bày hồ sơ HCCSVV xét theo ba hạng niên hạn và HCBVTQ xét theo thời gian phục vụ, kèm thanh tiến trình và gợi ý điều kiện cho từng cấp danh hiệu.

\begin{figure}[H]
    \centering
    \includegraphics[width=0.92\textwidth]{UI/07-admin-chi-tiet-quan-nhan-thong-tin}
    \caption{Chi tiết quân nhân, thẻ Thông tin cơ bản}
    \label{fig:ui-admin-personnel-info}
\end{figure}

\begin{figure}[H]
    \centering
    \includegraphics[width=0.92\textwidth]{UI/08-admin-chi-tiet-quan-nhan-ho-so-khen-thuong}
    \caption{Chi tiết quân nhân, thẻ Hồ sơ khen thưởng với thanh tiến trình điều kiện}
    \label{fig:ui-admin-personnel-award}
\end{figure}

\textbf{Duyệt đề xuất} gồm trang danh sách (Hình~\ref{fig:ui-admin-proposal-list}) với thẻ trạng thái và bộ lọc năm, loại, cùng cửa sổ chi tiết (Hình~\ref{fig:ui-admin-proposal-detail}) hiển thị tệp đính kèm cùng bảng danh hiệu của đề xuất kèm số quyết định trước khi phê duyệt hoặc từ chối.

\begin{figure}[H]
    \centering
    \includegraphics[width=0.92\textwidth]{UI/09-admin-duyet-de-xuat-danh-sach}
    \caption{Danh sách đề xuất chờ phê duyệt}
    \label{fig:ui-admin-proposal-list}
\end{figure}

\begin{figure}[H]
    \centering
    \includegraphics[width=0.92\textwidth]{UI/10-admin-duyet-de-xuat-chi-tiet}
    \caption{Chi tiết một đề xuất khen thưởng}
    \label{fig:ui-admin-proposal-detail}
\end{figure}

\textbf{Quản lý quyết định} gồm bảng danh sách quyết định (Hình~\ref{fig:ui-admin-decision-list}) với số QĐ, năm, ngày ký, người ký và thao tác (xem, sửa, xoá, xem PDF), cùng cửa sổ chi tiết (Hình~\ref{fig:ui-admin-decision-detail}) hiển thị đầy đủ thông tin và nút xem file PDF.

\begin{figure}[H]
    \centering
    \includegraphics[width=0.92\textwidth]{UI/11-admin-quan-ly-quyet-dinh}
    \caption{Quản lý danh sách quyết định khen thưởng}
    \label{fig:ui-admin-decision-list}
\end{figure}

\begin{figure}[H]
    \centering
    \includegraphics[width=0.85\textwidth]{UI/12-admin-chi-tiet-quyet-dinh}
    \caption{Chi tiết một quyết định khen thưởng}
    \label{fig:ui-admin-decision-detail}
\end{figure}

\textbf{Quản lý khen thưởng} (Hình~\ref{fig:ui-admin-award-list}) dùng hàng thẻ chọn loại danh hiệu, bộ lọc năm, họ tên, danh hiệu và bảng dữ liệu kèm số quyết định; cửa sổ \textbf{xuất Excel} (Hình~\ref{fig:ui-admin-award-export}) cho phép chọn khoảng thời gian, đơn vị, quân nhân và danh hiệu cần xuất.

\begin{figure}[H]
    \centering
    \includegraphics[width=0.92\textwidth]{UI/13-admin-quan-ly-khen-thuong}
    \caption{Quản lý khen thưởng theo loại danh hiệu}
    \label{fig:ui-admin-award-list}
\end{figure}

\begin{figure}[H]
    \centering
    \includegraphics[width=0.85\textwidth]{UI/14-admin-xuat-excel-khen-thuong}
    \caption{Cửa sổ xuất Excel danh sách khen thưởng}
    \label{fig:ui-admin-award-export}
\end{figure}

\textbf{Thêm khen thưởng đồng loạt} (Hình~\ref{fig:ui-admin-bulk}) áp dụng biểu mẫu sáu bước với stepper (chọn loại, chọn quân nhân, thiết lập danh hiệu, xem lại, thêm số quyết định, ghi nhận); tại bước chọn quân nhân, cán bộ có thể nhập nhanh từ tệp Excel hoặc chọn thủ công từ bảng.
\textbf{Khen thưởng đột xuất} (Hình~\ref{fig:ui-admin-adhoc}) quản lý các khen thưởng ngoài chu kỳ theo năm, đối tượng, hình thức và số quyết định.

\begin{figure}[H]
    \centering
    \includegraphics[width=0.92\textwidth]{UI/15-admin-them-khen-thuong-dong-loat}
    \caption{Thêm khen thưởng đồng loạt qua Excel}
    \label{fig:ui-admin-bulk}
\end{figure}

\begin{figure}[H]
    \centering
    \includegraphics[width=0.92\textwidth]{UI/16-admin-khen-thuong-dot-xuat}
    \caption{Quản lý khen thưởng đột xuất}
    \label{fig:ui-admin-adhoc}
\end{figure}

\textbf{Nhật ký hệ thống} (Hình~\ref{fig:ui-admin-logs}) hiện thực bố cục ``bộ lọc trên, bảng dưới'' đã mô tả ở wireframe (Mục~\ref{subsection:4.2.1}), với panel lọc nhanh theo thời gian, vai trò và hành động, bốn thẻ thống kê theo loại hành động và bảng nhật ký sắp xếp được theo cột.

\begin{figure}[H]
    \centering
    \includegraphics[width=0.92\textwidth]{UI/17-admin-nhat-ky-he-thong}
    \caption{Trang nhật ký hệ thống với bộ lọc và thống kê}
    \label{fig:ui-admin-logs}
\end{figure}

\subsubsection{Vai trò chỉ huy đơn vị (\texttt{Manager})}

\textbf{Trang chủ Chỉ huy đơn vị} (Hình~\ref{fig:ui-manager-dashboard}) giới hạn dữ liệu trong phạm vi đơn vị quản lý, với năm thẻ thống kê (quân số, tổng CSTĐCS, tổng NCKH, khen thưởng, đề xuất chờ duyệt) và các biểu đồ khen thưởng, đề xuất, NCKH.

\begin{figure}[H]
    \centering
    \includegraphics[width=0.92\textwidth]{UI/28-manager-dashboard}
    \caption{Trang chủ vai trò chỉ huy đơn vị}
    \label{fig:ui-manager-dashboard}
\end{figure}

\textbf{Đề xuất khen thưởng} (Hình~\ref{fig:ui-manager-proposal}) liệt kê các đề xuất đã gửi kèm hộp hướng dẫn, thẻ trạng thái và nút tạo đề xuất mới.
\textbf{Biểu mẫu tạo đề xuất} hiện thực wizard năm bước đã mô tả ở wireframe (Mục~\ref{subsection:4.2.1}); ba hình minh hoạ bước 1 chọn loại khen thưởng từ các thẻ bảy loại (Hình~\ref{fig:ui-manager-proposal-b1}), bước 2 chọn quân nhân từ bảng có lọc kèm cột hiển thị điều kiện (Hình~\ref{fig:ui-manager-proposal-b2}) và bước 5 xem lại toàn bộ thông tin trước khi gửi (Hình~\ref{fig:ui-manager-proposal-b5}).

\begin{figure}[H]
    \centering
    \includegraphics[width=0.92\textwidth]{UI/18-manager-de-xuat-khen-thuong}
    \caption{Danh sách đề xuất khen thưởng của đơn vị}
    \label{fig:ui-manager-proposal}
\end{figure}

\begin{figure}[H]
    \centering
    \includegraphics[width=0.92\textwidth]{UI/19-manager-tao-de-xuat-b1-chon-loai}
    \caption{Tạo đề xuất, bước 1: chọn loại khen thưởng}
    \label{fig:ui-manager-proposal-b1}
\end{figure}

\begin{figure}[H]
    \centering
    \includegraphics[width=0.92\textwidth]{UI/20-manager-tao-de-xuat-b2-chon-quan-nhan}
    \caption{Tạo đề xuất, bước 2: chọn quân nhân}
    \label{fig:ui-manager-proposal-b2}
\end{figure}

\begin{figure}[H]
    \centering
    \includegraphics[width=0.92\textwidth]{UI/30-manager-xem-lai-gui-de-xuat}
    \caption{Tạo đề xuất, bước 5: xem lại thông tin và gửi đề xuất}
    \label{fig:ui-manager-proposal-b5}
\end{figure}

\textbf{Lịch sử chức vụ} (Hình~\ref{fig:ui-manager-position}) ghi nhận các giai đoạn giữ chức vụ kèm hệ số và ba thẻ tổng thời gian theo nhóm hệ số.
\textbf{Danh hiệu hằng năm} của một quân nhân (Hình~\ref{fig:ui-manager-annual}) liệt kê các danh hiệu theo từng năm kèm số quyết định, còn \textbf{thành tích NCKH} (Hình~\ref{fig:ui-manager-nckh}) liệt kê các đề tài/sáng kiến theo năm.

\begin{figure}[H]
    \centering
    \includegraphics[width=0.92\textwidth]{UI/21-manager-lich-su-chuc-vu}
    \caption{Lịch sử chức vụ của quân nhân}
    \label{fig:ui-manager-position}
\end{figure}

\begin{figure}[H]
    \centering
    \includegraphics[width=0.92\textwidth]{UI/22-manager-khen-thuong-hang-nam}
    \caption{Danh hiệu khen thưởng hằng năm của quân nhân}
    \label{fig:ui-manager-annual}
\end{figure}

\begin{figure}[H]
    \centering
    \includegraphics[width=0.92\textwidth]{UI/23-manager-thanh-tich-nckh}
    \caption{Thành tích nghiên cứu khoa học của quân nhân}
    \label{fig:ui-manager-nckh}
\end{figure}

\textbf{Quản lý quân nhân đơn vị} (Hình~\ref{fig:ui-manager-personnel-list}) liệt kê quân số trong phạm vi đơn vị kèm ba thẻ thống kê và bộ lọc cấp bậc, chức vụ; từ mỗi dòng có thể mở chi tiết quân nhân.

\begin{figure}[H]
    \centering
    \includegraphics[width=0.92\textwidth]{UI/32-manager-quan-ly-quan-nhan-don-vi}
    \caption{Quản lý quân nhân thuộc đơn vị ở vai trò chỉ huy đơn vị}
    \label{fig:ui-manager-personnel-list}
\end{figure}

\textbf{Chi tiết quân nhân, thẻ Quản lý chi tiết} (Hình~\ref{fig:ui-manager-personnel-manage}) tập hợp các lối tắt quản lý theo lưới: lịch sử chức vụ, danh hiệu hằng năm, NCKH và các huân chương, huy chương, kỷ niệm chương.

\begin{figure}[H]
    \centering
    \includegraphics[width=0.92\textwidth]{UI/24-manager-chi-tiet-quan-nhan-quan-ly}
    \caption{Chi tiết quân nhân, thẻ Quản lý chi tiết}
    \label{fig:ui-manager-personnel-manage}
\end{figure}

\textbf{Thông báo giới hạn danh hiệu.} Khi xét một quân nhân đã đạt cấp danh hiệu cao nhất mà phần mềm hỗ trợ, hệ thống mở bảng tóm tắt hồ sơ kèm thông báo chưa hỗ trợ các danh hiệu cao hơn (Hình~\ref{fig:ui-chain-cap}), vì đây là danh hiệu chỉ trao một lần và không xét lại.

\begin{figure}[H]
    \centering
    \includegraphics[width=0.82\textwidth]{UI/29-phan-mem-chua-ho-tro}
    \caption{Thông báo giới hạn khi quân nhân đã đạt cấp danh hiệu cao nhất được hỗ trợ}
    \label{fig:ui-chain-cap}
\end{figure}

\subsubsection{Vai trò người dùng (\texttt{User})}

\textbf{Trang chủ Người dùng} (Hình~\ref{fig:ui-user-dashboard}) hiển thị hồ sơ cá nhân của quân nhân đăng nhập: các thẻ thống kê cá nhân (tổng CSTĐCS, NCKH, CSTĐCS liên tục, số tháng cống hiến), thông tin cá nhân và tóm tắt hồ sơ điều kiện hằng năm và niên hạn kèm gợi ý.
\textbf{Hồ sơ danh hiệu cá nhân hằng năm} (Hình~\ref{fig:ui-user-annual}) trình bày thanh tiến trình tổng hợp và dòng thời gian chuỗi danh hiệu theo năm (CSTĐCS $\rightarrow$ BKBQP $\rightarrow$ CSTĐTQ $\rightarrow$ BKTTCP).

\begin{figure}[H]
    \centering
    \includegraphics[width=0.92\textwidth]{UI/25-user-dashboard}
    \caption{Trang chủ vai trò người dùng}
    \label{fig:ui-user-dashboard}
\end{figure}

\begin{figure}[H]
    \centering
    \includegraphics[width=0.92\textwidth]{UI/26-user-ho-so-danh-hieu-hang-nam}
    \caption{Hồ sơ danh hiệu cá nhân hằng năm với dòng thời gian chuỗi danh hiệu}
    \label{fig:ui-user-annual}
\end{figure}

\subsection{Tóm tắt các thành phần kỹ thuật trọng yếu}
\label{subsection:4.3.4}

{\raggedright
Phần này tóm tắt năm thành phần kỹ thuật trọng yếu trong cách tổ chức nội bộ của hệ thống.

\textbf{Hàm tính ngữ cảnh chuỗi danh hiệu.} Từ danh sách danh hiệu hằng năm của một quân nhân cùng năm xét, thành phần này tính sẵn các chỉ số phục vụ việc xét chuỗi, gồm: năm bắt đầu chuỗi; năm gần nhất nhận mỗi cấp danh hiệu nếu vẫn nằm trong khoảng năm đang xét; số năm CSTĐCS liên tục kể từ lần nhận gần nhất; và số đợt đề nghị đã bỏ lỡ.
Vì là một hàm thuần, không truy vấn cơ sở dữ liệu, nó dễ kiểm chứng với nhiều bộ dữ liệu mẫu, từ chu kỳ vừa đủ đến lỡ nhiều chu kỳ.

\textbf{Chuỗi middleware bảo mật năm tầng.} Mỗi yêu cầu làm thay đổi dữ liệu lần lượt đi qua năm tầng: xác thực token, kiểm tra vai trò, kiểm tra dữ liệu bằng Zod, ghi nhật ký, rồi mới tới controller.
Việc kiểm tra vai trò và lọc dữ liệu theo phạm vi đơn vị không gom vào một middleware dùng chung. Mỗi route tự khai báo guard và bộ lọc cần thiết, nên chỉ áp đúng quy tắc của riêng route đó.
Một số route ghi hoặc tải tệp còn chèn thêm bước giới hạn tốc độ và bước nhận tệp đính kèm.
Vì mỗi tầng chỉ đảm nhiệm một việc nên dễ tái sử dụng giữa các route, đúng theo nguyên tắc đơn trách nhiệm.

\textbf{Tầng Repository tách khỏi Prisma Client.} Mỗi repository đóng gói các thao tác đọc, ghi dữ liệu sau những phương thức mang nghĩa nghiệp vụ, giấu đi chi tiết truy cập Prisma.
Mỗi phương thức có thể chạy bên trong một giao dịch chung khi được yêu cầu.
Nhờ đó tầng Service gom nhiều thao tác đọc ghi vào cùng một giao dịch, bảo đảm hoặc tất cả cùng thành công, hoặc cùng được hoàn tác khi có lỗi.
Mỗi chiến lược đề xuất gắn với tầng dữ liệu theo cùng nguyên tắc này: nó ghi dữ liệu qua \texttt{proposalRepository} bên trong một giao dịch, còn phần kiểm tra điều kiện được uỷ thác cho các hàm thuần trong gói \texttt{eligibility/}.

\textbf{Hàm xét điều kiện chuỗi.} Đây là thành phần cốt lõi của nghiệp vụ xét chuỗi danh hiệu.
Đầu vào gồm cấu hình của một cấp danh hiệu, ngữ cảnh chuỗi của quân nhân, cờ đánh dấu đã nhận hay chưa và số cờ tiền điều kiện nằm trong khung xét.
Hàm xử lý theo ba bước: chặn nhận lại nếu là danh hiệu một lần duy nhất mà quân nhân đã có; kiểm tra chuỗi đã đủ một chu kỳ trọn vẹn; và kiểm tra đủ số cờ tiền điều kiện trong khung xét cùng yêu cầu NCKH nếu cấp danh hiệu đòi hỏi.
Kết quả trả về là kết luận đủ hay chưa đủ điều kiện kèm thông điệp giải thích bằng tiếng Việt.

\textbf{Registry chiến lược cho bảy loại đề xuất.} Registry trung tâm ánh xạ tám loại khen thưởng tới chiến lược tương ứng.
Bảy loại đi qua luồng đề xuất và phê duyệt được gom vào registry này, mỗi loại ứng với một chiến lược tuân theo giao diện chung \texttt{ProposalStrategy}.
Riêng khen thưởng đột xuất (\texttt{DOT\_XUAT}) cố ý được ánh xạ tới \texttt{null} vì không đi qua luồng đề xuất và phê duyệt chuẩn.
Hai loại chỉ trao một huân chương duy nhất là HCQKQT và KNC dùng chung hàm \texttt{importSingleMedal} để tránh trùng lặp.
Nhờ cách tổ chức này, việc bổ sung một loại đề xuất mới chỉ cần viết thêm một chiến lược và đăng ký vào registry, không phải sửa tầng phê duyệt chung.\par}

\section{Kiểm thử}
\label{section:4.4}

Hệ thống được kiểm thử theo hai hướng bổ trợ nhau: kiểm thử hộp đen thủ công các chức năng nghiệp vụ qua giao diện (Mục~\ref{subsection:4.4.1}) và kiểm thử tự động bằng Jest cho các quy tắc nghiệp vụ lõi phía backend (Mục~\ref{subsection:4.4.2}).
Ngoài ra, sản phẩm còn được kiểm thử tương thích trên nhiều dòng máy và trình duyệt (Mục~\ref{subsection:4.4.3}).

\subsection{Kiểm thử hộp đen các chức năng nghiệp vụ}
\label{subsection:4.4.1}

Các chức năng nghiệp vụ chính được kiểm thử hộp đen thủ công thông qua giao diện trên môi trường phát triển.
Tám bảng kiểm thử dưới đây tóm tắt kết quả.

\begin{table}[H]
\centering
\caption{Kiểm thử chức năng đăng nhập}
\label{tab:kt-dangnhap}
\setlength{\tabcolsep}{4pt}\renewcommand{\arraystretch}{1.25}\begin{tabular}{|c|R{3.2cm}|R{3.3cm}|R{4.5cm}|c|}
\hline
\textbf{STT} & \textbf{Chức năng} & \textbf{Đầu vào} & \textbf{Đầu ra mong muốn} & \textbf{Kết quả} \\ \hline
1 & Đăng nhập hợp lệ & Tên đăng nhập và mật khẩu đúng & Chuyển hướng tới trang chủ quản trị viên & Đạt \\ \hline
2 & Sai mật khẩu & Tên đăng nhập đúng, mật khẩu sai & Thông báo chung, không tiết lộ trường nào sai & Đạt \\ \hline
3 & Tài khoản không tồn tại & Tên đăng nhập không có trong hệ thống & Thông báo chung, không lộ thông tin & Đạt \\ \hline
4 & Vượt giới hạn đăng nhập sai & Gửi quá số lần đăng nhập sai cho phép & Tạm thời từ chối yêu cầu đăng nhập & Đạt \\ \hline
5 & Làm mới phiên đăng nhập & Phiên đăng nhập hết hạn & Hệ thống tự động cấp phiên mới & Đạt \\ \hline
\end{tabular}
\end{table}

\begin{table}[H]
\centering
\caption{Kiểm thử chức năng quản lý quân nhân}
\label{tab:kt-quannhan}
\setlength{\tabcolsep}{4pt}\renewcommand{\arraystretch}{1.25}\begin{tabular}{|c|R{3.2cm}|R{3.3cm}|R{4.5cm}|c|}
\hline
\textbf{STT} & \textbf{Chức năng} & \textbf{Đầu vào} & \textbf{Đầu ra mong muốn} & \textbf{Kết quả} \\ \hline
1 & Thêm mới hợp lệ & Nhập đầy đủ thông tin & Bản ghi mới xuất hiện trong danh sách & Đạt \\ \hline
2 & Trùng CCCD & Nhập CCCD đã tồn tại & Hiển thị lỗi ``CCCD đã tồn tại trong hệ thống'' & Đạt \\ \hline
3 & Thiếu trường bắt buộc & Bỏ trống họ tên & Biểu mẫu chặn gửi và báo lỗi & Đạt \\ \hline
4 & Sửa thông tin & Thay đổi cấp bậc & Lưu thành công & Đạt \\ \hline
5 & Chuyển đơn vị & Chọn đơn vị mới & Lịch sử khen thưởng giữ nguyên, số lượng quân nhân của đơn vị được cập nhật & Đạt \\ \hline
\end{tabular}
\end{table}

\begin{table}[H]
\centering
\caption{Kiểm thử chức năng đề xuất khen thưởng}
\label{tab:kt-dexuat}
\setlength{\tabcolsep}{4pt}\renewcommand{\arraystretch}{1.25}\begin{tabular}{|c|R{3.2cm}|R{3.3cm}|R{4.5cm}|c|}
\hline
\textbf{STT} & \textbf{Chức năng} & \textbf{Đầu vào} & \textbf{Đầu ra mong muốn} & \textbf{Kết quả} \\ \hline
1 & Tạo đề xuất hợp lệ & Chọn 5 quân nhân đủ điều kiện & Đề xuất ở trạng thái chờ duyệt & Đạt \\ \hline
2 & Giới hạn theo đơn vị phụ trách & Chỉ huy đơn vị mở danh sách chọn quân nhân & Danh sách chỉ hiển thị quân nhân thuộc đơn vị phụ trách, không có quân nhân đơn vị khác & Đạt \\ \hline
3 & Quân nhân chưa đủ điều kiện & Đề nghị BKBQP cho quân nhân mới 1 năm CSTĐCS & Danh hiệu chưa đủ điều kiện bị loại khỏi danh sách chọn; nếu cố gửi qua API, backend từ chối kèm lý do & Đạt \\ \hline
4 & Lưu nháp khi tải lại trang & Tải lại trang giữa chừng & Dữ liệu nhập dở được giữ nguyên và tự xoá sau khi gửi hoặc đóng trang & Đạt \\ \hline
\end{tabular}
\end{table}

\begin{table}[H]
\centering
\caption{Kiểm thử chức năng phê duyệt đề xuất}
\label{tab:kt-pheduyet}
\setlength{\tabcolsep}{4pt}\renewcommand{\arraystretch}{1.25}\begin{tabular}{|c|R{3.2cm}|R{3.3cm}|R{4.5cm}|c|}
\hline
\textbf{STT} & \textbf{Chức năng} & \textbf{Đầu vào} & \textbf{Đầu ra mong muốn} & \textbf{Kết quả} \\ \hline
1 & Phê duyệt thông thường & Đề xuất hợp lệ & Đề xuất chuyển trạng thái đã duyệt và lưu vào hồ sơ khen thưởng & Đạt \\ \hline
2 & Phê duyệt kèm gắn quyết định & Số quyết định, tháng và năm nhận & Đề xuất được gắn quyết định, hồ sơ quân nhân hiển thị đúng dữ liệu & Đạt \\ \hline
3 & Từ chối đề xuất & Bấm ``Từ chối'' kèm lý do & Đề xuất chuyển trạng thái bị từ chối, gửi thông báo cho Chỉ huy đơn vị & Đạt \\ \hline
4 & Đề nghị trùng BKTTCP & Đề nghị BKTTCP cho người đã có & Hệ thống chặn và huỷ bỏ thao tác & Đạt \\ \hline
5 & Hai cán bộ cùng phê duyệt & Hai phiên đồng thời & Chỉ một thành công, người sau bị từ chối do xung đột & Đạt \\ \hline
\end{tabular}
\end{table}

\begin{table}[H]
\centering
\caption{Kiểm thử chức năng tính lại điều kiện khen thưởng hằng năm}
\label{tab:kt-tinhlai}
\setlength{\tabcolsep}{4pt}\renewcommand{\arraystretch}{1.25}\begin{tabular}{|c|R{3.2cm}|R{3.3cm}|R{4.5cm}|c|}
\hline
\textbf{STT} & \textbf{Chức năng} & \textbf{Đầu vào} & \textbf{Đầu ra mong muốn} & \textbf{Kết quả} \\ \hline
1 & Quân nhân chưa nhận & 2 năm CSTĐCS liên tục & Hệ thống đề nghị BKBQP kèm lời giải thích & Đạt \\ \hline
2 & Đã nhận BKTTCP & Quân nhân đã có BKTTCP & Hiển thị thông báo đã có BKTTCP và chưa hỗ trợ danh hiệu cao hơn & Đạt \\ \hline
3 & Bỏ lỡ đợt đề nghị & 6 năm CSTĐCS liên tục, bỏ lỡ mốc năm 2 và năm 4 & Vẫn đủ điều kiện đề nghị BKBQP vì chu kỳ lặp mỗi 2 năm & Đạt \\ \hline
4 & Cửa sổ xét 3 năm của CSTĐTQ & BKBQP đã nhận cách 4 năm & BKBQP đó không được tính, yêu cầu BKBQP mới & Đạt \\ \hline
5 & Thiếu nghiên cứu khoa học & 3 năm CSTĐCS nhưng thiếu một năm NCKH & Chưa đủ điều kiện, gợi ý cần bổ sung nghiên cứu khoa học & Đạt \\ \hline
\end{tabular}
\end{table}

\begin{table}[H]
\centering
\caption[Kiểm thử nhập Excel hàng loạt]{Kiểm thử chức năng nhập Excel hàng loạt cho danh hiệu hằng năm}
\label{tab:kt-excel}
\setlength{\tabcolsep}{4pt}\renewcommand{\arraystretch}{1.25}\begin{tabular}{|c|R{3.2cm}|R{3.3cm}|R{4.5cm}|c|}
\hline
\textbf{STT} & \textbf{Chức năng} & \textbf{Đầu vào} & \textbf{Đầu ra mong muốn} & \textbf{Kết quả} \\ \hline
1 & Nhập 50 bản ghi hợp lệ & Tệp Excel đúng mẫu & Tạo đủ 50 bản ghi và ghi nhật ký hệ thống & Đạt \\ \hline
2 & Có dòng trùng dữ liệu & Một dòng trùng cùng quân nhân và năm & Bản xem trước báo 49 dòng hợp lệ và 1 dòng lỗi & Đạt \\ \hline
3 & Dữ liệu thay đổi giữa hai bước & Có bản ghi chen vào sau khi xem trước & Huỷ bỏ toàn bộ, không ghi bản ghi nào & Đạt \\ \hline
4 & Tệp sai định dạng & Tải lên tệp không phải Excel & Hệ thống báo lỗi định dạng tệp & Đạt \\ \hline
\end{tabular}
\end{table}

\begin{table}[H]
\centering
\caption{Kiểm thử chức năng quản lý đơn vị và chức vụ}
\label{tab:kt-donvi}
\setlength{\tabcolsep}{4pt}\renewcommand{\arraystretch}{1.25}\begin{tabular}{|c|R{3.2cm}|R{3.3cm}|R{4.5cm}|c|}
\hline
\textbf{STT} & \textbf{Chức năng} & \textbf{Đầu vào} & \textbf{Đầu ra mong muốn} & \textbf{Kết quả} \\ \hline
1 & Thêm cơ quan đơn vị mới & Nhập tên và mã đơn vị hợp lệ & Đơn vị mới xuất hiện trong cây tổ chức & Đạt \\ \hline
2 & Xoá đơn vị trực thuộc có quân nhân & Chọn đơn vị trực thuộc vẫn còn quân nhân & Hệ thống từ chối xoá và báo lỗi ràng buộc dữ liệu & Đạt \\ \hline
3 & Sửa hệ số chức vụ & Thay đổi hệ số từ 0{,}8 sang 0{,}9 & Lưu hệ số mới thành công; áp dụng cho bản ghi lịch sử chức vụ tạo sau đó & Đạt \\ \hline
\end{tabular}
\end{table}

\begin{table}[H]
\centering
\caption{Kiểm thử chức năng quản trị (sao lưu, nhật ký)}
\label{tab:kt-quantri}
\setlength{\tabcolsep}{4pt}\renewcommand{\arraystretch}{1.25}\begin{tabular}{|c|R{3.2cm}|R{3.3cm}|R{4.5cm}|c|}
\hline
\textbf{STT} & \textbf{Chức năng} & \textbf{Đầu vào} & \textbf{Đầu ra mong muốn} & \textbf{Kết quả} \\ \hline
1 & Sao lưu thủ công & Chọn ``Sao lưu ngay'' & Tạo tệp sao lưu mới & Đạt \\ \hline
2 & Phân quyền xem nhật ký sao lưu & Chỉ huy đơn vị mở trang nhật ký & Không hiển thị các bản ghi nhật ký sao lưu & Đạt \\ \hline
3 & Đổi lịch sao lưu & Đổi chu kỳ từ 24 giờ sang 12 giờ & Hệ thống tự sao lưu theo lịch mới & Đạt \\ \hline
\end{tabular}
\end{table}

\subsection{Kiểm thử tự động bằng Jest}
\label{subsection:4.4.2}

Bên cạnh kiểm thử hộp đen thủ công, các quy tắc nghiệp vụ lõi phía backend được bao quát bởi một bộ kiểm thử tự động viết bằng Jest, tập trung vào những quy tắc dễ sai và khó kiểm soát bằng tay.
Bộ kiểm thử gồm tổng cộng 82 tệp, phân theo nhóm chức năng được tóm tắt ở Bảng~\ref{tab:kt-jest}.
Phần lớn là kiểm thử đơn vị chạy trên Prisma được giả lập nên độc lập với cơ sở dữ liệu; riêng nhóm tích hợp gồm một tệp chạy trên cơ sở dữ liệu thật để kiểm chứng các thao tác có giao dịch.

\begin{table}[H]
\centering
\caption{Phân nhóm các tệp kiểm thử Jest}
\label{tab:kt-jest}
\renewcommand{\arraystretch}{1.25}
\begin{tabular}{|R{7.5cm}|c|}
\hline
\textbf{Nhóm kiểm thử} & \textbf{Số tệp} \\ \hline
Dịch vụ nghiệp vụ & 34 \\ \hline
Kịch bản chuỗi danh hiệu và đua tranh & 12 \\ \hline
Phê duyệt đề xuất & 10 \\ \hline
Nhập Excel & 7 \\ \hline
Gửi đề xuất & 7 \\ \hline
Phân quyền & 5 \\ \hline
Xác thực phiên đăng nhập & 3 \\ \hline
Hằng số & 2 \\ \hline
Xác thực dữ liệu đầu vào & 1 \\ \hline
Kiểm thử tích hợp & 1 \\ \hline
\textbf{Tổng} & \textbf{82} \\ \hline
\end{tabular}
\end{table}

Trọng tâm là quy tắc xét chuỗi danh hiệu, với các kịch bản: chu kỳ vừa đủ, chu kỳ thừa, lỡ một và lỡ nhiều chu kỳ, cờ tiền điều kiện rơi khỏi cửa sổ trượt, và việc chặn nhận lại danh hiệu một lần duy nhất.
Ngoài ra, bộ kiểm thử còn phủ luồng phê duyệt đề xuất, phân quyền theo bốn vai trò, cơ chế xoay vòng refresh token và việc xác thực dữ liệu đầu vào.
Nhờ tách phần quy tắc thành các hàm thuần (Mục~\ref{subsection:4.1.4}), những quy tắc này được kiểm thử độc lập với cơ sở dữ liệu, qua đó bảo đảm hai luồng xét điều kiện, dù khi tính lại hồ sơ hay khi phê duyệt, đều cho kết quả nhất quán.

\subsection{Kiểm thử tương thích}
\label{subsection:4.4.3}

Hệ thống được kiểm thử trên năm dòng máy phổ biến với các tổ hợp hệ điều hành và trình duyệt khác nhau nhằm đánh giá khả năng hiển thị và vận hành đồng nhất.

\begin{table}[H]
\centering
\caption{Kiểm thử tương thích các dòng máy}
\label{tab:kt-trinhduyet}
\renewcommand{\arraystretch}{1.3}
\begin{tabular}{|R{3.2cm}|R{4.2cm}|R{3.6cm}|R{2.4cm}|}
\hline
\textbf{Thiết bị} & \textbf{Phần cứng} & \textbf{Phần mềm} & \textbf{Kết quả} \\ \hline
MacBook Pro M1 Pro 2021 & Apple M1 Pro, RAM 16 GB, màn hình 14,2 inch Liquid Retina XDR & macOS Sonoma 14, Safari 17 & Đạt\textsuperscript{*} \\ \hline
Lenovo Legion 5 2022 & AMD Ryzen 7 5800H, RAM 16 GB, màn hình 15,6 inch FHD 165 Hz & Windows 11, Microsoft Edge 119 & Đạt \\ \hline
Lenovo IdeaPad Gaming 3 & AMD Ryzen 5 5600H, RAM 8 GB, màn hình 15,6 inch FHD & Windows 11, Google Chrome 120 & Đạt \\ \hline
Dell Inspiron 15 3520 & Intel Core i5-1235U, RAM 16 GB, màn hình 15,6 inch FHD & Windows 10, Mozilla Firefox 121 & Đạt \\ \hline
Mac mini M4 2024 & Apple M4, RAM 16 GB, màn hình ngoài 23,8 inch FHD & macOS Sequoia 15, Google Chrome 120 & Đạt \\ \hline
\end{tabular}

\smallskip
{\footnotesize\raggedright \textsuperscript{*}Trên trình duyệt Safari, bộ chọn ngày có khác biệt nhỏ về hiển thị nhưng không ảnh hưởng đến chức năng chính.\par}
\end{table}

\section{Triển khai}
\label{section:4.5}

\subsection{Yêu cầu phần cứng và phần mềm}
\label{subsection:4.5.1}

Hệ thống được thiết kế để triển khai trên một máy chủ Windows Server đơn lẻ, yêu cầu cấu hình phần cứng và phần mềm tối thiểu như sau:

\begin{itemize}
    \item CPU: 2 nhân, 2{,}0 GHz trở lên.
    \item RAM: 4 GB (khuyến nghị 8 GB cho môi trường vài chục người dùng).
    \item Ổ cứng: 50 GB SSD cho dữ liệu và bản sao lưu.
    \item Hệ điều hành: Windows Server 2019 hoặc Windows Server 2022.
    \item Phần mềm phụ thuộc: Node.js 20 LTS, PostgreSQL 15.
\end{itemize}

\subsection{Hướng dẫn triển khai môi trường phát triển}
\label{subsection:4.5.2}

Các bước cài đặt trên máy phát triển cá nhân được trình bày tuần tự dưới đây.

\textbf{Bước 1.} Cài đặt phụ thuộc Node bằng lệnh \texttt{npm install} lần lượt trong hai thư mục \texttt{BE-QLKT} và \texttt{FE-QLKT}.
Lệnh này tải các gói phụ thuộc liệt kê tại tệp \texttt{package.json} của mỗi module.

\textbf{Bước 2.} Sao chép tệp cấu hình mẫu \texttt{.env.example} thành \texttt{.env} trong cả hai thư mục, rồi chỉnh các giá trị thiết yếu: \texttt{DATABASE\_URL} cho chuỗi kết nối PostgreSQL, \texttt{JWT\_SECRET} cho khoá phát hành Access Token, \texttt{JWT\_REFRESH\_SECRET} cho khoá phát hành Refresh Token, \texttt{CLIENT\_URL} cho địa chỉ frontend.

\textbf{Bước 3.} Khởi tạo cơ sở dữ liệu trong thư mục \texttt{BE-QLKT} bằng lệnh \texttt{npx prisma db push} để đồng bộ lược đồ Prisma vào cơ sở dữ liệu; tiếp theo chạy lệnh \texttt{npm run init-super-admin} để tạo tài khoản Quản trị viên mặc định.

\textbf{Bước 4.} Khởi chạy môi trường phát triển bằng \texttt{npm run dev} ở \texttt{BE-QLKT} cho backend tại cổng 4000 và lệnh tương tự ở \texttt{FE-QLKT} cho frontend tại cổng 3000.

Sau bước 4, mở trình duyệt tới \url{http://localhost:3000} và đăng nhập bằng tài khoản Quản trị viên mặc định để bắt đầu cấu hình các đơn vị, chức vụ và tài khoản phụ trợ.

\subsection{Triển khai sản xuất với PM2}
\label{subsection:4.5.3}

Trên máy chủ sản xuất, sản phẩm được vận hành bằng PM2, một công cụ quản lý tiến trình hỗ trợ tự khởi động lại khi tiến trình bị dừng đột ngột, ghi log có thời gian và theo dõi tài nguyên hệ thống.
Tệp \texttt{ecosystem.config.js} ở thư mục gốc của dự án khai báo hai tiến trình ứng dụng: \texttt{be-qlkt} chạy bản đã build (\texttt{dist/index.js}) trong thư mục \texttt{BE-QLKT} ở cổng 4000 và \texttt{fe-qlkt} chạy \texttt{next start} trong thư mục \texttt{FE-QLKT} ở cổng 3000.
Cả hai tiến trình được đặt biến môi trường \texttt{NODE\_ENV} là \texttt{production}.

Khởi chạy bằng lệnh \texttt{pm2 start ecosystem.config.js}, sau đó dùng \texttt{pm2 save} để ghi nhớ danh sách tiến trình đang chạy.
Để PM2 tự khởi động cùng máy chủ Windows, cài gói \texttt{pm2-windows-startup} và chạy \texttt{pm2-startup install}, đăng ký PM2 như một tác vụ khởi động của Windows.
Người dùng trong mạng nội bộ truy cập giao diện do tiến trình frontend phục vụ ở cổng 3000; trình duyệt gọi trực tiếp tới backend ở cổng 4000 qua REST API và Socket.IO.

\subsection{Sao lưu và giám sát}
\label{subsection:4.5.4}

Tác vụ sao lưu chạy bên trong tiến trình backend nhờ \texttt{node-cron}; lịch chạy và trạng thái được cấu hình ngay trong DevZone.
Mặc định, hệ thống sao lưu lúc 02:00 hằng ngày: mỗi lần chạy, dịch vụ đọc 21 bảng dữ liệu nghiệp vụ (không gồm nhật ký hệ thống và thông báo), kết xuất thành các câu lệnh chèn dữ liệu SQL đã làm sạch rồi ghi ra một tệp \texttt{.sql} trong thư mục sao lưu.
Các bản cũ quá số ngày lưu trữ đã cấu hình (mặc định 15 ngày) được tự động dọn.
Về giám sát, PM2 cung cấp lệnh xem log thời gian thực và gom toàn bộ log vào một thư mục riêng, có thể tích hợp với hệ thống giám sát tập trung của Học viện trong tương lai.

\section*{Kết chương}

Chương này đã trình bày kiến trúc phân tầng, thiết kế gói, thiết kế lớp, thiết kế cơ sở dữ liệu, thiết kế giao diện, kết quả xây dựng, kết quả kiểm thử hộp đen các chức năng nghiệp vụ và phương án triển khai của hệ thống.
Các kết quả này cho thấy phần mềm đã hiện thực hoá đầy đủ những yêu cầu chức năng và phi chức năng đặt ra ở Chương~\ref{chapter:Related_works}.
Chương tiếp theo đi sâu phân tích những đóng góp nổi bật của đồ án.

\end{document}
